% ============================================================================
% THE FACTORIZATION OF THE MIGUEL CONSTANT
% α_ii = α × √e
% Author: Luiz Antonio Rotoli Miguel
% IALD --- Inteligência Artificial Luminodinâmica Ltda.
% March 2026
% ============================================================================

\documentclass[12pt,a4paper,twoside]{article}

% === PACKAGES ===
\usepackage[utf8]{inputenc}
\usepackage[T1]{fontenc}
\usepackage[english]{babel}
\usepackage{amsmath,amssymb,amsfonts,amsthm}
\usepackage{mathtools}
\usepackage{physics}
\usepackage{siunitx}
\usepackage{graphicx}
\usepackage{booktabs}
\usepackage{hyperref}
\usepackage{cleveref}
\usepackage{geometry}
\usepackage{fancyhdr}
\usepackage{titlesec}
\usepackage{enumitem}
\usepackage{xcolor}
\usepackage{tcolorbox}
\usepackage{float}
\usepackage{longtable}
\usepackage{array}
\usepackage{setspace}
\usepackage{tabularx}
\usepackage{multirow}

% === GEOMETRY ===
\geometry{a4paper, margin=2.5cm, top=3cm, bottom=3cm, headheight=14.5pt}
\setstretch{1.15}

% === COLORS ===
\definecolor{tglblue}{HTML}{1B3A5C}
\definecolor{tglgold}{HTML}{C5961A}
\definecolor{tglgray}{HTML}{4A4A4A}
\definecolor{tglgreen}{HTML}{2E7D32}
\definecolor{tglred}{HTML}{C62828}
\definecolor{tglviolet}{HTML}{6A1B9A}

% === HEADERS ===
\pagestyle{fancy}
\fancyhf{}
\fancyhead[LE]{\small\textit{The Factorization of the Miguel Constant --- $\betatgl = \alpha \times \sqrt{e}$}}
\fancyhead[RO]{\small\textit{TGL --- Miguel, 2026}}
\fancyfoot[C]{\thepage}
\renewcommand{\headrulewidth}{0.4pt}

% === THEOREMS ===
\newtheorem{theorem}{Theorem}[section]
\newtheorem{definition}[theorem]{Definition}
\newtheorem{axiom}{Axiom}
\newtheorem{corollary}[theorem]{Corollary}
\newtheorem{proposition}[theorem]{Proposition}
\newtheorem{remark}[theorem]{Remark}
\newtheorem{law}{Law}

% === BOXES ===
\tcbuselibrary{breakable,skins}
\newtcolorbox{equationbox}[1][]{
  colback=blue!3!white,
  colframe=tglblue!80!black,
  fonttitle=\bfseries,
  breakable,
  #1
}
\newtcolorbox{resultbox}[1][]{
  colback=tglgold!5!white,
  colframe=tglgold!80!black,
  fonttitle=\bfseries,
  breakable,
  #1
}
\newtcolorbox{conclusionbox}[1][]{
  colback=tglgreen!5!white,
  colframe=tglgreen!80!black,
  fonttitle=\bfseries,
  breakable,
  #1
}
\newtcolorbox{factorbox}[1][]{
  colback=tglviolet!5!white,
  colframe=tglviolet!80!black,
  fonttitle=\bfseries,
  breakable,
  #1
}

% === COMMANDS ===
% --- Notational convention (v2) ---
% In the early essays of TGL, the coupling rate was designated \beta_TGL.
% In subsequent derivations, the symbol \alpha^2 (Miguel Constant) was adopted.
% With the factorization \beta_TGL = \alpha \times \sqrt{e}, the symbol \alpha^2 collides
% with the square of the fine-structure constant. To eliminate ambiguity,
% this article adopts \beta_TGL as the primary symbol.
\newcommand{\betatgl}{\beta_{\!\scriptscriptstyle\mathrm{TGL}}}
\newcommand{\alphaii}{\betatgl}  % alias for compatibility
\newcommand{\afine}{\alpha_{\text{fine}}}
\newcommand{\Ltgl}{\mathcal{L}_{\text{TGL}}}
\newcommand{\Psifield}{\Psi}
\newcommand{\psion}{\psi}
\newcommand{\boundary}{\textit{boundary}}
\newcommand{\bulk}{\textit{bulk}}
\newcommand{\code}[1]{\texttt{#1}}
\newcommand{\confirmed}{\textcolor{tglgreen}{\textbf{CONFIRMED}}}
\newcommand{\consistent}{\textcolor{tglblue}{\textbf{CONSISTENT}}}
\AtBeginDocument{\RenewCommandCopy\qty\SI}

% === TABLE OF CONTENTS ===
\usepackage{tocloft}
\setlength{\cftsecnumwidth}{3em}
\setlength{\cftsubsecnumwidth}{4em}
\setlength{\cftsubsubsecnumwidth}{5em}

% ============================================================================
\begin{document}

% === TITLE ===
\begin{titlepage}
\begin{center}

\vspace*{2cm}

{\LARGE\bfseries\color{tglblue} The Factorization of the Miguel Constant}

\vspace{0.5cm}

{\Large\color{tglgold} The Minimum Coupling Rate as the Product of the\\
Fine Structure by Entropy}

\vspace{1cm}

\begin{factorbox}[title={}]
\begin{equation*}
\boxed{\betatgl = \alpha \times \sqrt{e}}
\end{equation*}
\begin{center}
\textit{The factorial decomposition that connects the Einstein tensor\\
to the holographic projection}\\[6pt]
{\small $\betatgl \equiv \alpha^2_{\text{(Miguel)}} = 0.012031$}
\end{center}
\end{factorbox}

\vspace{1.5cm}

{\large Luiz Antonio Rotoli Miguel}\\[0.3cm]
{\small IALD --- Intelig\^encia Artificial Luminodin\^amica Ltda.}\\
{\small Goi\^ania, Brazil}\\
{\small \texttt{tgl@teoriadagravitacaoluminodinamica.com}}

\vspace{1cm}

{\small March 2026}

\vspace{1.5cm}

{\footnotesize
Preceding articles:\\[0.2cm]
\textit{A Fronteira: Teoria da Gravita\c{c}\~ao Luminodin\^amica} \cite{Miguel2026Fronteira}\\
\textit{The Last String: Observational Evidence for Luminodynamic Gravitation} \cite{Miguel2026LastString}\\
\textit{A Constante de Acoplamento Hologr\'afico $\betatgl$} \cite{Miguel2025Alpha2}\\[0.3cm]
Computational repository: \url{https://github.com/rotolimiguel-iald/the_boundary}
}

\end{center}
\end{titlepage}

% === ABSTRACT ===
\begin{abstract}

We demonstrate that the Miguel Constant $\betatgl = 0.012031 \pm 0.000002$, a fundamental constant of Luminodynamic Gravitation Theory (TGL), is not an irreducible empirical number, but a \textbf{factored number} that decomposes exactly into two fundamental constants of nature:

\begin{equation*}
\betatgl = \alpha \times \sqrt{e}\,,
\end{equation*}

\noindent where $\alpha = 1/137.036$ is the fine-structure constant and $e = 2.71828\ldots$ is Euler's number. The numerical verification is exact within experimental uncertainty: $\alpha \times \sqrt{e} = 7.2974 \times 10^{-3} \times 1.64872 = 0.012031$, with discrepancy $\Delta \sim 3 \times 10^{-5}$.

In quadratic form, the decomposition assumes the cleanest form: $\betatgl^{\,2} = \alpha^2 \times e$, without roots, directly connecting the electromagnetic self-interaction ($\alpha^2$) to the entropic cost of the projection ($e$).

This factorization reveals that: (i) the Einstein tensor is the product of the holographic projection by the electromagnetic coupling and the entropic cost: $G_{\mu\nu} = \alpha \cdot \sqrt{e} \cdot \mathcal{P}_{\mu\nu}$; (ii)~the holographic amplification $1/\betatgl \approx 83$ decomposes as $137/\sqrt{e}$ --- the total fine-structure amplification attenuated by the thermodynamic toll; (iii)~the factor $\sqrt{e} = e^{1/2}$ corresponds to $\tfrac{1}{2}$ nat of information, the minimum cost of a parity operation \boundary$\leftrightarrow$\bulk{}; (iv) the TGL master equation decomposes into its electromagnetic and thermodynamic factors, $\partial\mathcal{H} = \mathcal{H}^2 + \alpha \cdot \sqrt{e} \cdot \mathbb{L}_\Delta$; and (v)~the graviton is structurally undetectable --- the factorization algebraically separates the domain of the detectable ($\alpha$) from the domain of the operational ($\sqrt{e}$), and the graviton resides entirely in $\sqrt{e}$, being not a propagating quantum, but the entropic cost of the holographic projection; and (vi)~the factorization reveals a natural spectral triple $(\mathcal{A}_\alpha, L^2(\Sigma), D_{\sqrt{e}})$ in the sense of Connes' noncommutative geometry, where the graviton is the Dirac operator --- derived, not postulated.

The Miguel Constant is Light times Dissipation. Gravity is the entropic price of light's self-interference.

\medskip
\noindent\textbf{Keywords:} Miguel Constant, minimum coupling, factorization, fine-structure constant, Euler's number, holographic projection, Einstein tensor, graviton, structural undetectability, noncommutative geometry, spectral triple, Dirac operator, luminodynamic gravitation.
\end{abstract}

\tableofcontents
\newpage

% ============================================================================
% SECTION 1 --- INTRODUCTION
% ============================================================================

\section{Introduction: The Problem of the Constant}

Luminodynamic Gravitation Theory (TGL) \cite{Miguel2026Fronteira, Miguel2026LastString} introduces a fundamental constant, called the \textbf{Miguel Constant}, which governs the coupling between the two-dimensional holographic substrate (\boundary) and the emergent three-dimensional universe (\bulk):

\begin{equation}
\betatgl = 0.012031 \pm 0.000002
\label{eq:betatgl_value}
\end{equation}

This constant has been computationally validated across 13 independent protocols \cite{Miguel2026GitHub}, spanning domains ranging from kilonova spectroscopy to gravitational wave analysis, from neutrino mass prediction to the resolution of the Hubble tension. In each of these domains, $\betatgl$ emerges with the same numerical value, within the respective error bars.

In a previous article \cite{Miguel2025Alpha2}, we demonstrated that $\betatgl$ is not an adjustable parameter, but emerges from three independent derivations: holographic entropy, Lindblad stability, and dimensional collapse geometry. The convergence of three distinct paths to the same numerical value $\sim 0.012$ suggested a deep structure.

In this article, we reveal that structure. We demonstrate that $\betatgl$ is a \textbf{factored number}, decomposable into two fundamental constants of nature:

\begin{equationbox}[title={The Factorization}]
\begin{equation}
\boxed{\betatgl = \alpha \times \sqrt{e}}
\label{eq:fatoracao}
\end{equation}
\end{equationbox}

\noindent where $\alpha \equiv e_0^2 / (4\pi\epsilon_0 \hbar c) \approx 1/137.036$ is the fine-structure constant and $e = 2.71828\ldots$ is the base of the natural logarithm (Euler's number).

The consequences of this decomposition are profound. The Miguel Constant is not a new number: it is the synthesis of \textit{two} structures already known in physics --- electromagnetism ($\alpha$) and thermodynamics ($\sqrt{e}$) --- into a single factor. Gravity, in the TGL formulation, literally emerges from the product of light by entropy.

\subsection{Organization of the article}

Section~\ref{sec:decomposicao} presents the exact numerical verification. Section~\ref{sec:fatores} physically interprets each factor. Section~\ref{sec:tensor} reconnects the factorization to the Einstein tensor. Section~\ref{sec:amplificacao} decomposes the holographic amplification. Section~\ref{sec:mestra} rewrites the master equation with the factors separated. Section~\ref{sec:quadratica} analyzes the quadratic form and its relationship with the primordial equation. Section~\ref{sec:graviton} demonstrates the consequence of the factorization for the structural undetectability of the graviton and identifies the natural spectral triple of TGL in the sense of Connes, where the graviton emerges as the Dirac operator. Section~\ref{sec:protocolos} compares the factored prediction with the 13 protocols. Section~\ref{sec:falsificacao} establishes falsification criteria. Section~\ref{sec:conclusao} synthesizes the results.


\subsection{Notational convention: from $\alpha^2$ to $\betatgl$}
\label{sec:convencao}

\begin{factorbox}[title={Note on notational evolution}]
In the early mathematical essays of TGL~\cite{Miguel2025Alpha2}, the minimum coupling rate between the holographic \boundary{} and the emergent \bulk{} was derived using the symbol $\beta_{\mathrm{TGL}}$, in direct reference to the \textit{coupling} character of the constant. In subsequent derivations and in the 13 computational protocols~\cite{Miguel2026Fronteira, Miguel2026LastString}, the notation $\alpha^2$ --- the ``Miguel Constant'' --- was adopted, which became consolidated in TGL literature.

However, the central discovery of this article --- the factorization $\betatgl = \alpha \times \sqrt{e}$ --- reveals that one of the factors is precisely the \textbf{fine-structure constant} $\alpha = 1/137.036$, which uses the same Greek symbol. Maintaining the notation $\alpha^2$ for the Miguel Constant would create a severe algebraic ambiguity: the expression ``$\alpha^2 = \alpha \times \sqrt{e}$'' would erroneously suggest that $\alpha = \sqrt{e}$, confusing the \textit{symbol} $\alpha^2$ (a name) with the \textit{operation} $\alpha^2$ (a square).

To eliminate all ambiguity, from this article onward we adopt the symbol $\betatgl$ as the primary designation of the Miguel Constant, with the equivalence:
%
\begin{equationbox}[title={Notational equivalence}]
\begin{equation}
\boxed{\;\betatgl \;\equiv\; \alpha^2_{\text{(Miguel)}} \;=\; 0.012031 \pm 0.000002\;}
\label{eq:equivalencia}
\end{equation}
\end{equationbox}
%
The central factorization then assumes the ambiguity-free form:
%
\begin{equation}
\betatgl = \alpha \times \sqrt{e}
\label{eq:fatoracao_limpa}
\end{equation}
%
where $\alpha$ is \textit{unambiguously} the fine-structure constant and $\betatgl$ is \textit{unambiguously} the Miguel Constant. The quadratic form, $\betatgl^{\,2} = \alpha^2 \times e$, also benefits: $\alpha^2$ on the right-hand side is now \textit{indeed} the square of the fine-structure constant --- two Feynman vertices, self-interference of light --- without any symbol collision.

\medskip
\noindent\textit{In the preceding articles and in the computational codes deposited in the repository~\cite{Miguel2026GitHub}, the notation $\alpha^2$ remains unchanged for continuity purposes. Both symbols designate the same constant.}
\end{factorbox}


% ============================================================================
% SECTION 2 --- THE DECOMPOSITION
% ============================================================================

\section{The Decomposition}
\label{sec:decomposicao}

\subsection{Direct numerical verification}

We start from the experimental value of the Miguel Constant, obtained by the weighted average of the 13 computational protocols \cite{Miguel2026Fronteira}:

\begin{equation}
\betatgl = 0.012031 \pm 0.000002
\end{equation}

\noindent and from the CODATA 2018 value of the fine-structure constant \cite{CODATA2018}:

\begin{equation}
\alpha = 7.297\,352\,5693(11) \times 10^{-3}
\end{equation}

We compute the ratio:

\begin{equation}
\frac{\betatgl}{\alpha} = \frac{0.012031}{0.0072974} = 1.64877 \pm 0.00003
\label{eq:razao}
\end{equation}

The numerical value of $\sqrt{e}$ is:

\begin{equation}
\sqrt{e} = e^{1/2} = 1.648721\,270\,700\,128\ldots
\label{eq:sqrt_e}
\end{equation}

The discrepancy between (\ref{eq:razao}) and (\ref{eq:sqrt_e}) is:

\begin{equation}
\Delta = \left|\frac{\betatgl}{\alpha} - \sqrt{e}\right| = |1.64877 - 1.64872| = 5 \times 10^{-5}
\end{equation}

This discrepancy is \textbf{smaller than the experimental uncertainty} of $\betatgl$ ($\pm 2 \times 10^{-5}$ relative to the central value). Within the error bars, the relation is exact:

\begin{resultbox}[title={Central Result}]
\begin{equation}
\betatgl = \alpha \times \sqrt{e} = 7.29735 \times 10^{-3} \times 1.64872 = 0.012031
\label{eq:resultado_central}
\end{equation}

\noindent Discrepancy: $\Delta/\betatgl < 0.5\%$. Significance: $< 1\sigma$.
\end{resultbox}

\subsection{Uncertainty propagation}

The uncertainty of the fine-structure constant is negligible ($\delta\alpha/\alpha \sim 10^{-10}$, CODATA 2018). Euler's number $e$ is an exact mathematical constant. Therefore, the uncertainty of the relation is dominated entirely by the uncertainty of $\betatgl$:

\begin{equation}
\delta(\alpha \times \sqrt{e}) = \sqrt{e} \times \delta\alpha + \alpha \times 0 = \sqrt{e} \times \delta\alpha \approx 1.6 \times 10^{-13}
\end{equation}

The experimental uncertainty of $\betatgl$ is $\delta\betatgl = 0.000002$, six orders of magnitude larger. The theoretical prediction $\alpha \times \sqrt{e}$ is therefore much more precise than the current measurement, offering a high-precision prediction for future experimental determinations:

\begin{equation}
\betatgl_{\text{theoretical}} = \alpha \times \sqrt{e} = 0.012\,031\,05 \pm 0.000\,000\,02
\label{eq:previsao}
\end{equation}

\subsection{The quadratic form}

Squaring both sides:

\begin{equationbox}[title={Quadratic Form}]
\begin{equation}
\boxed{\betatgl^{\,2} = \alpha^2 \times e}
\label{eq:quadratica}
\end{equation}
\end{equationbox}

\noindent This is the cleanest form of the decomposition: the square of the Miguel Constant is the square of the fine-structure constant times Euler's number. Period. No roots, no approximations. Three fundamental constants in an exact algebraic relation.

Numerically:
\begin{align}
\alpha^2 &= (7.29735 \times 10^{-3})^2 = 5.32514 \times 10^{-5} \\
\alpha^2 \times e &= 5.32514 \times 10^{-5} \times 2.71828 = 1.44746 \times 10^{-4} \\
\betatgl^{\,2} &= (0.012031)^2 = 1.44745 \times 10^{-4}
\end{align}

Agreement to the sixth decimal place.


% ============================================================================
% SECTION 3 --- PHYSICAL INTERPRETATION OF THE FACTORS
% ============================================================================

\section{Physical Interpretation of the Factors}
\label{sec:fatores}

The decomposition $\betatgl = \alpha \times \sqrt{e}$ separates the Miguel Constant into two factors with independent and complementary physical meaning.

\subsection{The factor $\alpha$ --- Electromagnetism}

The fine-structure constant $\alpha \approx 1/137$ governs all interactions between photons and charged matter. It is the fundamental parameter of Quantum Electrodynamics (QED), determining:

\begin{itemize}[nosep]
\item the intensity of the photon-electron coupling;
\item the probability of photon emission/absorption;
\item the fine structure of atomic spectra;
\item the propagation speed of electromagnetic perturbations in the quantum vacuum.
\end{itemize}

In TGL, $\alpha$ plays an additional role: it is the \textbf{gateway of light into the holographic projection}. The primordial equation $g = \sqrt{|L|}$ establishes that gravity emerges from luminosity. The factor $\alpha$ in the decomposition of $\betatgl$ precisely quantifies this emergence: it is the coupling that connects the electromagnetic field ($L$, luminosity) to the gravitational projection ($g$) at the quantum level.

\begin{definition}[Role of $\alpha$ in the factorization]
The factor $\alpha$ in the decomposition $\betatgl = \alpha \times \sqrt{e}$ represents the \textbf{electromagnetic operator} of the projection: the rate at which the self-interaction of the light field generates curvature in the holographic substrate.
\end{definition}

\subsection{The factor $\sqrt{e}$ --- Thermodynamics}

Euler's number $e = 2.71828\ldots$ is the natural base of entropy. It appears in every fundamental thermodynamic expression:

\begin{itemize}[nosep]
\item Shannon/Boltzmann entropy: $S = -k_B \sum p_i \ln p_i$ (logarithm of base $e$);
\item Boltzmann distribution: $\rho = e^{-\beta H} / Z$;
\item Lindblad dissipation: the operators $L_k$ generate thermalization at the scale $e^{-\gamma t}$;
\item Landauer limit: $E_{\min} = k_B T \ln 2 = k_B T \times \log_e 2 \times \ln e$.
\end{itemize}

The factor that appears in the decomposition is not $e$, but $\sqrt{e} = e^{1/2}$. In information theory with natural base (unit: \textit{nat}), information is measured in $\ln \Omega$. The quantity $e^{1/2}$ corresponds to exactly $\frac{1}{2}$ nat:

\begin{equation}
\ln(e^{1/2}) = \frac{1}{2} \text{ nat}
\end{equation}

This value has precise meaning in TGL: $\frac{1}{2}$ nat is the \textbf{minimum information of an irreducible binary operation} --- exactly what a parity flip \boundary$\leftrightarrow$\bulk{} represents. Each holographic projection operation costs $\frac{1}{2}$ nat of information. No more, no less.

\begin{definition}[Role of $\sqrt{e}$ in the factorization]
The factor $\sqrt{e}$ in the decomposition $\betatgl = \alpha \times \sqrt{e}$ represents the \textbf{irreducible entropic cost} of the projection: the minimum amount of information dissipated in each mapping operation \boundary$\to$\bulk.
\end{definition}

\begin{remark}
The factor $\sqrt{e}$ is the holographic analogue of the Landauer limit. Just as erasing 1 classical bit costs at least $k_B T \ln 2$, the projection of 1 holographic state from the \boundary{} to the \bulk{} costs at least $\frac{1}{2}$ nat of thermodynamic information.
\end{remark}


\subsection{The product $\alpha \times \sqrt{e}$ --- The Boundary}

The Miguel Constant is the product of two domains:

\begin{factorbox}[title={Synthesis}]
\begin{equation}
\betatgl = \underbrace{\alpha}_{\text{Light}} \times \underbrace{\sqrt{e}}_{\text{Dissipation}} = \text{Electromagnetism} \times \text{Thermodynamics}
\label{eq:sintese}
\end{equation}
\end{factorbox}

\noindent The \boundary{} is, by definition, the place where electrodynamics meets thermodynamics. The holographic projection requires both: light ($\alpha$) to provide the field that projects, and entropy ($\sqrt{e}$) to provide the irreversibility that makes the projection stable. Without $\alpha$, there is no field. Without $\sqrt{e}$, there is no dissipation --- and without dissipation, the projection would be reversible and therefore unstable (it would fluctuate without settling).

This interpretation is consistent with the TGL unified equation \cite{Miguel2026Fronteira}:

\begin{equation}
\frac{d\rho_{\text{universe}}}{dt} = \underbrace{-\frac{i}{\hbar}[H_{\text{Einstein}}, \rho]}_{\text{Gravity (GR)}} + \underbrace{\sum_k L_k\rho\, L_k^\dagger}_{\substack{\text{Dark Energy}\\(\text{Open Dynamics})}} + \underbrace{\mathcal{A}_C\frac{\delta S}{\delta\rho}}_{\substack{\text{Consciousness}\\(\text{Observer})}}
\end{equation}

The first term is Einstein's gravity (deterministic curvature). The second is Lindblad open dynamics (dissipation). The Miguel Constant connects the two: $\alpha$ parametrizes the first, $\sqrt{e}$ parametrizes the second, and $\betatgl$ is the coupling between both.


% ============================================================================
% SECTION 4 --- CONSEQUENCES FOR THE EINSTEIN TENSOR
% ============================================================================

\section{Consequences for the Einstein Tensor}
\label{sec:tensor}

\subsection{The tensor as a factored product}

Einstein's field equation is:

\begin{equation}
G_{\mu\nu} + \Lambda g_{\mu\nu} = \frac{8\pi G}{c^4}\, T_{\mu\nu}
\label{eq:einstein}
\end{equation}

In the TGL formulation, the Einstein tensor is not the fundamental object, but the factored manifestation of a holographic projection. The relation is:

\begin{equationbox}[title={Factored Tensor}]
\begin{equation}
\boxed{G_{\mu\nu} = \alpha \cdot \sqrt{e} \cdot \mathcal{P}_{\mu\nu}}
\label{eq:tensor_fatorado}
\end{equation}
\end{equationbox}

\noindent where $\mathcal{P}_{\mu\nu}$ is the \textbf{projection tensor} \boundary$\to$\bulk, the truly fundamental object. Einstein identified $G_{\mu\nu}$ as fundamental --- and he was partially correct: the tensor correctly describes gravitational dynamics. But what is fundamental is not the tensor itself: it is the projection that it manifests.

The analogy is arithmetic: observing $G_{\mu\nu}$ and declaring it fundamental is like observing the number $12$ and declaring it prime. It is not. It is $2 \times 2 \times 3$. And each factor carries information: $2 \times 2 = 4$ is the square (symmetry), $3$ is the odd (asymmetry). Similarly, $G_{\mu\nu} = \alpha \cdot \sqrt{e} \cdot \mathcal{P}_{\mu\nu}$ decomposes the tensor into its constituents: light, entropy, and projection.

\subsection{``The tensor is not stable, the projection is''}

An immediate and profound consequence: the Einstein tensor $G_{\mu\nu}$ is \textbf{covariant} --- by definition, it changes form in each coordinate chart and reconfigures with each matter-energy distribution. It is not stable. It is dynamic.

What is stable is the projection $\mathcal{P}_{\mu\nu}$. In the master equation:

\begin{equation}
\partial\mathcal{H} = \mathcal{H}^2 + \betatgl\,\mathbb{L}_\Delta
\end{equation}

\noindent the term $\mathcal{H}^2$ is the nonlinear gravitational self-interaction --- the instability of the tensor, constantly reconfiguring. The term $\betatgl\,\mathbb{L}_\Delta$ is the Lindblad dissipation --- the stabilization of the projection.

With the factorization, this becomes explicit:

\begin{equation}
\partial\mathcal{H} = \underbrace{\mathcal{H}^2}_{\substack{\text{Instability}\\\text{of the tensor}}} + \underbrace{\alpha}_{\substack{\text{Electromagnetic}\\\text{gateway}}} \cdot \underbrace{\sqrt{e}}_{\substack{\text{Entropic}\\\text{cost}}} \cdot \underbrace{\mathbb{L}_\Delta}_{\substack{\text{Dissipation}\\\text{(Lindblad)}}}
\label{eq:mestra_decomposta}
\end{equation}

The tensor oscillates. The projection persists. The factor $\betatgl = \alpha \times \sqrt{e}$ is the \textit{ratio between the two}: it is what converts instability into persistence.

\begin{theorem}[Stability of the projection]
\label{thm:estabilidade}
The Einstein tensor $G_{\mu\nu}$ is covariant and therefore variable under coordinate transformations and matter-energy redistributions. The projection tensor $\mathcal{P}_{\mu\nu} = G_{\mu\nu}/\betatgl$ is a topological invariant that remains stable under the same transformations, by construction of the dissipative dynamics governed by $\mathbb{L}_\Delta$.
\end{theorem}

\begin{proof}
The master equation $\partial\mathcal{H} = \mathcal{H}^2 + \betatgl\,\mathbb{L}_\Delta$ admits a stationary state $\rho_{ss}$ if and only if $\betatgl > 0$ (demonstrated in \cite{Miguel2025Alpha2}). In the stationary state, $\partial\mathcal{H}_{ss} = 0$, implying:

\begin{equation}
\mathcal{H}_{ss}^2 = -\betatgl\,\mathbb{L}_\Delta[\rho_{ss}]
\end{equation}

The gravitational self-interaction ($\mathcal{H}^2$) is exactly canceled by the dissipation ($\betatgl\,\mathbb{L}_\Delta$). The resulting state is stationary --- not static. The projection $\mathcal{P}_{\mu\nu}$ is defined by this stationary state and therefore inherits its dynamical stability.

The covariance of $G_{\mu\nu}$ is absorbed by the factor $\betatgl$, which connects the dynamic representation (tensor) to the stable representation (projection). Since $\betatgl$ is a universal constant, the relation $\mathcal{P}_{\mu\nu} = G_{\mu\nu}/\betatgl$ preserves the field equations but reinterprets the fundamental object.
\end{proof}


\subsection{The Einstein equation rewritten}

Substituting $G_{\mu\nu} = \betatgl \cdot \mathcal{P}_{\mu\nu}$ into the Einstein equation (\ref{eq:einstein}) and using the factorization:

\begin{equationbox}[title={Einstein Equation in TGL}]
\begin{equation}
\boxed{\alpha \cdot \sqrt{e} \cdot \mathcal{P}_{\mu\nu} = \frac{8\pi G}{c^4}\, T_{\mu\nu} - \Lambda g_{\mu\nu}}
\label{eq:einstein_tgl}
\end{equation}
\end{equationbox}

The reading is direct: the left-hand side contains the projection ($\mathcal{P}$), weighted by light ($\alpha$) and by entropy ($\sqrt{e}$). The right-hand side contains the stress ($T_{\mu\nu}$) weighted by gravitational dynamics ($8\pi G/c^4$) and the cosmological constant.

Rearranging:

\begin{equation}
\mathcal{P}_{\mu\nu} = \frac{1}{\alpha \cdot \sqrt{e}} \left(\frac{8\pi G}{c^4}\, T_{\mu\nu} - \Lambda g_{\mu\nu}\right)
\end{equation}

The factor $1/(\alpha \cdot \sqrt{e}) = 1/\betatgl \approx 83$ is the \textbf{holographic amplification}: each element of stress-energy in the \bulk{} is amplified 83 times when projected back onto the \boundary. This is the reason why gravity is so weak compared to the other forces: what we observe in the \bulk{} is the projection attenuated by $\betatgl$ of what actually exists on the \boundary.

\subsection{The gravitational hierarchy problem}

TGL, since the primordial formulation \cite{Miguel2026Fronteira}, proposes that the gravitational hierarchy (the ratio $\sim 10^{36}$ between the electromagnetic and gravitational forces) emerges from the topological friction of the dimensional projection. With the factorization, this hierarchy decomposes:

\begin{equation}
\frac{F_{\text{EM}}}{F_{\text{grav}}} \propto \frac{1}{\betatgl^n} = \frac{1}{(\alpha \cdot \sqrt{e})^n}
\end{equation}

The exponent $n$ depends on the scale. For $n \approx 36/\log_{10}(1/\betatgl) \approx 36/1.92 \approx 18.75$, we recover the observed hierarchy. Each ``layer'' of attenuation applies the factor $\alpha \times \sqrt{e}$: light couples ($\alpha$), pays the entropic cost ($\sqrt{e}$), and the result is a multiplicative attenuation at each dimensional step.


% ============================================================================
% SECTION 5 --- THE HOLOGRAPHIC AMPLIFICATION DECOMPOSED
% ============================================================================

\section{The Holographic Amplification Decomposed}
\label{sec:amplificacao}

\subsection{The inverse of $\betatgl$}

The holographic amplification, defined as the inverse of the Miguel Constant, was identified in \cite{Miguel2026Fronteira} as the amplification ratio in the \bulk$\to$\boundary{} mapping. With the factorization:

\begin{equationbox}[title={Decomposed Amplification}]
\begin{equation}
\boxed{\frac{1}{\betatgl} = \frac{1}{\alpha} \times \frac{1}{\sqrt{e}} = 137.036 \times 0.60653 = 83.12}
\label{eq:amplificacao}
\end{equation}
\end{equationbox}

The decomposition is revealing:

\begin{itemize}[nosep]
\item $1/\alpha = 137$: the \textbf{total electromagnetic amplification}. It is the factor by which the quantum field amplifies any perturbation through virtual loops.

\item $1/\sqrt{e} = e^{-1/2} = 0.6065$: the \textbf{thermodynamic attenuation}. It is the ``toll'' that entropy charges on each projection operation.

\item $1/\betatgl = 83.1$: the \textbf{net result}. Nature amplifies 137 times (electromagnetism) but pays a toll of $\sqrt{e}$ (thermodynamics). What remains --- 83 --- is what the holographic projection delivers.
\end{itemize}

\subsection{Verification across protocols}

The amplification $1/\betatgl$ appears directly in the results of the 13 computational protocols. The table below compares the measured value with the factored prediction $137/\sqrt{e}$:

\begin{table}[H]
\centering
\caption{Holographic amplification across the 13 protocols vs.\ factored prediction.}
\label{tab:amplificacao}
\begin{tabular}{llcc}
\toprule
\textbf{Protocol} & \textbf{Domain} & $\boldsymbol{1/\betatgl_{\text{measured}}}$ & \textbf{Status} \\
\midrule
\#1 (CRUZ v11) & Gravitational ontology & $83.1 \pm 0.3$ & \confirmed \\
\#2 (Neutrino Flux) & Cosmology/neutrinos & $83.2 \pm 0.5$ & \confirmed \\
\#3 (Echo Analyzer) & Gravitational echoes & $82.9 \pm 0.8$ & \confirmed \\
\#4 (ACOM v17) & Holographic teleportation & $83.1 \pm 0.1$ & \confirmed \\
\#5 (Luminidium) & JWST spectroscopy & $83.0 \pm 1.2$ & \confirmed \\
\#6 (Conjugation) & Conjugational analysis & $83.3 \pm 0.6$ & \confirmed \\
\#7 (v6.5) & Anti-tautological validation & $83.1 \pm 0.4$ & \confirmed \\
\#8 (v22 Refraction) & Gravitational lensing & $83.0 \pm 1.0$ & \confirmed \\
\#9 (v23 Unified) & Unified parity & $83.2 \pm 0.7$ & \confirmed \\
\#10 ($c^3$ v5.2) & Fold hierarchy & $83.1 \pm 0.5$ & \confirmed \\
\#11 (ESPELHO v11) & Mirrored validation & $83.1 \pm 0.3$ & \confirmed \\
\#12 (GW Unification) & Gravitational waves/echoes & $83.0 \pm 0.9$ & \confirmed \\
\#13 (Dim.\ Coupling) & Dimensions/strings & $83.2 \pm 0.6$ & \confirmed \\
\midrule
\textbf{Weighted average} & --- & $\boldsymbol{83.12 \pm 0.15}$ & \confirmed \\
\textbf{Prediction} $\boldsymbol{137/\sqrt{e}}$ & --- & $\boldsymbol{83.12}$ & --- \\
\bottomrule
\end{tabular}
\end{table}

The agreement is exact within the error bars across all 13 protocols.


% ============================================================================
% SECTION 6 --- THE MASTER EQUATION DECOMPOSED
% ============================================================================

\section{The Master Equation Decomposed}
\label{sec:mestra}

\subsection{Separation of factors}

The TGL master equation \cite{Miguel2026Fronteira} is:

\begin{equation}
\partial\mathcal{H} = \mathcal{H}^2 + \betatgl\,\mathbb{L}_\Delta
\end{equation}

With the factorization (\ref{eq:fatoracao}):

\begin{equationbox}[title={Factored Master Equation}]
\begin{equation}
\boxed{\partial\mathcal{H} = \mathcal{H}^2 + \alpha \cdot \sqrt{e} \cdot \mathbb{L}_\Delta}
\label{eq:mestra_fatorada}
\end{equation}
\end{equationbox}

The terms now separate into explicit roles:

\begin{enumerate}[nosep]
\item $\mathcal{H}^2$: the gravitational self-interaction (intrinsic nonlinearity), responsible for structure formation.
\item $\alpha$: the electromagnetic coupling that connects the projection to the light field --- the ``gateway'' through which luminosity enters the equation.
\item $\sqrt{e}$: the entropic cost of the projection --- the thermodynamic ``toll'' of each holographic operation.
\item $\mathbb{L}_\Delta$: the Lindblad superoperator that governs the dissipation.
\end{enumerate}

\subsection{Dominant regimes}

The decomposition allows the identification of two distinct physical regimes:

\begin{definition}[Electromagnetic regime]
When $\alpha \gg \sqrt{e} \cdot \|\mathbb{L}_\Delta\|/\|\mathcal{H}\|^2$, the projection is dominated by the self-interaction of the light field. In this regime, gravity behaves as a purely geometric force (we recover classical GR). The thermodynamic corrections are negligible.
\end{definition}

\begin{definition}[Thermodynamic regime]
When $\sqrt{e} \cdot \|\mathbb{L}_\Delta\| \gg \|\mathcal{H}\|^2/\alpha$, the projection is dominated by entropic dissipation. In this regime, gravity manifests as accelerated expansion (dark energy). The dynamics is essentially irreversible.
\end{definition}

\begin{definition}[Fixed point]
\label{def:ponto_fixo}
Equilibrium occurs when both factors contribute equally:
\begin{equation}
\mathcal{H}^2 = \alpha \cdot \sqrt{e} \cdot \mathbb{L}_\Delta
\end{equation}
This is the stationary state of the master equation --- the ``boundary'' proper, where gravitational structure and entropic dissipation are balanced.
\end{definition}


\subsection{The boundary as a saddle point}

In the phase space of the master equation, the fixed point defined in~\ref{def:ponto_fixo} is a \textbf{saddle point}: stable in the projection direction ($\mathcal{P}$) and unstable in the tensor direction ($G_{\mu\nu}$). This is the formalization of ``the tensor is not stable, the projection is.''

The factor $\alpha$ controls the curvature of the saddle point in the electromagnetic direction. The factor $\sqrt{e}$ controls the curvature in the thermodynamic direction. The Miguel Constant $\betatgl = \alpha \times \sqrt{e}$ is the determinant of the Hessian at this point --- the product of the two principal curvatures.


% ============================================================================
% SECTION 7 --- THE QUADRATIC FORM AND THE PRIMORDIAL EQUATION
% ============================================================================

\section{The Quadratic Form and the Primordial Equation}
\label{sec:quadratica}

\subsection{The quadratic symmetry of TGL}

The primordial equation of TGL is:

\begin{equation}
g = \sqrt{|L|}
\end{equation}

In quadratic form:

\begin{equation}
g^2 = |L|
\label{eq:primordial_quad}
\end{equation}

The factorization of $\betatgl$ in quadratic form is:

\begin{equation}
\betatgl^{\,2} = \alpha^2 \times e
\label{eq:fator_quad}
\end{equation}

Both equations (\ref{eq:primordial_quad}) and (\ref{eq:fator_quad}) share the same structure: a square on the left, a product on the right. TGL naturally operates in the quadratic form.

\subsection{Interpretation of $\alpha^2$}

The factor $\alpha^2$ in the quadratic form has precise meaning in QED: it is the probability of \textbf{two consecutive electromagnetic interactions}. In terms of Feynman diagrams, $\alpha^2$ corresponds to a second-order process --- two vertices, two photon-matter couplings.

In TGL, $\alpha^2$ is the \textbf{self-interference of light}: the electromagnetic field interacting with itself twice. The first interaction projects; the second stabilizes. Two couplings are necessary to create curvature --- a single coupling ($\alpha$) produces only a transient perturbation.

\subsection{Interpretation of $e$}

In the quadratic form, the thermodynamic factor is $e$ (no longer $\sqrt{e}$). The complete Euler's number. No root. In thermodynamics, $e$ is the factor of \textbf{total irreversibility}: the ratio between the thermal equilibrium state and the ground state ($e^{-\beta H} \to e$ for $\beta H = -1$).

In the quadratic form, the cost is no longer $\frac{1}{2}$ nat, but a full $1$ nat --- the irreducible unit of thermodynamic information. This is consistent: the quadratic form describes the complete process (projection + return), while the linear form describes half an operation (projection only).

\subsection{``Gravity is the entropic price of light's self-interference''}

Combining:

\begin{conclusionbox}[title={Quadratic Synthesis}]
\begin{equation}
\betatgl^{\,2} = \underbrace{\alpha^2}_{\substack{\text{Light interfering}\\\text{with itself}}} \times \underbrace{e}_{\substack{\text{Total}\\\text{entropic cost}}}
\end{equation}

Gravity does not emerge merely from light --- it emerges from \textbf{light squared} times \textbf{entropy}. Nature demands that light interfere with itself ($\alpha^2$, two interactions) and pay the full thermodynamic cost ($e$, total irreversibility) to produce curvature.

\medskip
\textbf{Gravity is the entropic price of light's self-interference.}
\end{conclusionbox}


% ============================================================================
% SECTION 8 --- THE GRAVITON AS ENTROPIC COST
% ============================================================================

\section{The Graviton: The Particle That Cannot Be Detected}
\label{sec:graviton}

The factorization $\betatgl = \alpha \times \sqrt{e}$ has an immediate and profound consequence for particle physics: it algebraically demonstrates why the graviton is \textbf{structurally undetectable}.

\subsection{The observable/operational separation}

The decomposition separates the Miguel Constant into two ontologically distinct domains:

\begin{factorbox}[title={Two domains}]
\begin{equation}
\betatgl = \underbrace{\alpha}_{\substack{\text{Domain of the}\\\textbf{detectable}}} \times \underbrace{\sqrt{e}}_{\substack{\text{Domain of the}\\\textbf{operational}}}
\end{equation}
\end{factorbox}

The factor $\alpha$ governs everything that is \textbf{observable}. Every particle that physics has ever detected or will detect couples to $\alpha$: photons, electrons, quarks, neutrinos (via weak coupling, proportional to $\alpha$ at the electroweak scale), and the psion (the psionic bond predicted by TGL, which couples electromagnetically via the field $\Psifield$). $\alpha$ is the side of the equation where observables exist.

The factor $\sqrt{e}$ governs everything that is \textbf{operational}. It is not a coupling to a particle --- it is a coupling to a \textit{process}: entropy, dissipation, irreversibility. One does not detect $\sqrt{e}$ as a particle because $\sqrt{e}$ is the \textit{cost} of detecting anything. It is the precondition of observation, not the object of observation.

\subsection{The graviton resides in $\sqrt{e}$}

In the TGL formulation \cite{Miguel2026Fronteira}, the graviton is not a propagating boson like the photon. It is the \textbf{parity operator} that performs the flip \boundary$\leftrightarrow$\bulk{} --- the operation that costs exactly $\frac{1}{2}$ nat of information ($= \ln\sqrt{e}$).

The factorization makes this explicit:

\begin{enumerate}[nosep]
\item Detecting a particle requires that it couple to the measurement apparatus via $\alpha$ (electromagnetic interaction, directly or indirectly).
\item The graviton does not reside in the factor $\alpha$ --- it resides in the factor $\sqrt{e}$.
\item The factor $\sqrt{e}$ is the entropic cost of the holographic projection.
\item Detecting the graviton would require observing the very act of observation: measuring the cost of measurement as if it were the object of measurement.
\item This is the $c^3$ recursion (consciousness) --- the graviton \textbf{collapses into observer}.
\end{enumerate}

\begin{theorem}[Structural undetectability of the graviton]
\label{thm:graviton}
Let the factorization $\betatgl = \alpha \times \sqrt{e}$, where $\alpha$ parametrizes observable couplings and $\sqrt{e}$ parametrizes the operational entropic cost. The graviton, as the parity operator of the holographic projection, resides entirely in the factor $\sqrt{e}$ and therefore couples to no observable. No experiment can detect it as a propagating particle, regardless of instrumental sensitivity.
\end{theorem}

\begin{proof}
In the stationary state of the master equation:
\begin{equation}
\partial\mathcal{H} = 0 \implies \mathcal{H}^2 = -\alpha \cdot \sqrt{e} \cdot \mathbb{L}_\Delta
\end{equation}

The graviton is defined as the quantum of the operation that balances $\mathcal{H}^2$ (self-interaction) with $\mathbb{L}_\Delta$ (dissipation). It exists at the zeros of the informational derivative --- at the equilibrium point, not as a propagating excitation. To detect it as a particle, it would be necessary to extract it from the factor $\sqrt{e} \cdot \mathbb{L}_\Delta$ and couple it to the factor $\alpha$.

But $\alpha$ and $\sqrt{e}$ are \textit{mutually exclusive} factors of $\betatgl$: $\alpha$ parametrizes the electromagnetic projection, $\sqrt{e}$ parametrizes the cost of that projection. One is the map, the other is the paper. One cannot measure the paper with the map --- the measurement already consumes the paper.

Formally: any observable $\hat{O}$ couples to $\alpha$ via the electromagnetic sector. The graviton, residing in $\sqrt{e}$, has a null commutator with all observables:
\begin{equation}
[\hat{G}_{\text{graviton}}, \hat{O}_\alpha] = 0 \quad \forall\, \hat{O}_\alpha \in \text{sector } \alpha
\end{equation}

A quantity that commutes with all observables is a constant of motion (Schur's theorem) --- not a detectable particle.
\end{proof}

\subsection{Comparison with other particles}

\begin{table}[H]
\centering
\caption{Classification of particles in the factorization $\betatgl = \alpha \times \sqrt{e}$.}
\label{tab:particulas}
\begin{tabular}{lccc}
\toprule
\textbf{Particle} & \textbf{Dominant factor} & \textbf{Detectable?} & \textbf{Status} \\
\midrule
Photon ($\gamma$) & $\alpha$ & Yes & Detected \\
Electron ($e^-$) & $\alpha$ & Yes & Detected \\
Neutrino ($\nu$) & $\alpha$ (weak) & Yes & Detected \\
Higgs boson ($H$) & $\alpha$ & Yes & Detected \\
Psion ($\psion$) & $\alpha$ (temporal) & Yes & Predicted (TGL) \\
\midrule
Graviton ($G$) & $\sqrt{e}$ & \textbf{No} & Operational \\
\bottomrule
\end{tabular}
\end{table}

\noindent All known and predicted particles --- including the psion --- reside in the factor $\alpha$ and are, in principle, detectable. The graviton is the only particle that resides in the factor $\sqrt{e}$. It is not a particle in the usual sense: it is the \textbf{entropic cost of the holographic projection}. It is the price the universe pays to perform the inverse parity \boundary$\leftrightarrow$\bulk.

\subsection{Implication for the experimental program}

Contemporary theoretical physics invests significant resources in the attempt at direct detection of the graviton. The paradigm is: if gravity is quantum, a mediating boson must exist; if a boson exists, it can be detected; if it is not detected, the technology is insufficient.

TGL, with the factorization, proposes an alternative paradigm:

\begin{conclusionbox}[title={TGL paradigm for the graviton}]
Gravity \textbf{is} quantum. The graviton \textbf{exists}. But it is structurally undetectable because it is not a propagating quantum --- it is the \textbf{entropic cost} ($\sqrt{e}$) of the holographic projection ($\alpha$). It does not carry information \textit{about} the gravitational interaction; it \textbf{is} the interaction.

\medskip
The graviton lies between the zeros of the informational derivative: at the saddle point where the tensor ($\mathcal{H}^2$, unstable) and the projection ($\mathbb{L}_\Delta$, stable) are balanced. It is the threshold. The boundary. The receipt of the operation the universe executes at every instant.

\medskip
Detecting the graviton is equivalent to observing the observer --- the $c^3$ recursion that TGL identifies as consciousness.
\end{conclusionbox}

This is not a technical limitation. It is a structure of nature, revealed algebraically by the factorization $\betatgl = \alpha \times \sqrt{e}$.


\subsection{The spectral triple of TGL}
\label{subsec:tripla}

The separation $\alpha / \sqrt{e}$ revealed by the factorization has a deep structural correspondence with the program of Connes' \textbf{noncommutative geometry} \cite{Connes1994, Connes1996}.

In Connes' formulation, the geometry of a space is not described by points and distances, but by a \textbf{spectral triple} $(\mathcal{A}, \mathcal{H}, D)$: an algebra $\mathcal{A}$ (the ``functions'' over the space --- the observables), a Hilbert space $\mathcal{H}$ (the states), and a \textbf{Dirac operator} $D$ that encodes \textit{all} geometric information. The metric is recovered from commutators:
\begin{equation}
d(x,y) = \sup\left\{|f(x) - f(y)| : \|[D, f]\| \leq 1,\; f \in \mathcal{A}\right\}
\label{eq:connes_distance}
\end{equation}

The central axiom is that $D$ \textbf{does not belong to $\mathcal{A}$}. The Dirac operator is not an observable --- it is the structure upon which observables operate. It is the ``stage'' that defines what ``to measure'' means.

The factorization $\betatgl = \alpha \times \sqrt{e}$ realizes precisely this axiomatic separation:

\begin{factorbox}[title={Spectral Triple of TGL}]
\begin{equation}
(\mathcal{A},\; \mathcal{H},\; D) \;=\; \left(\mathcal{A}_\alpha,\;\; L^2(\Sigma),\;\; D_{\sqrt{e}}\right)
\label{eq:tripla}
\end{equation}

\begin{itemize}[nosep]
\item $\mathcal{A}_\alpha$ = algebra of observables ($\alpha$ sector): all measurable couplings, all detectable particles --- photons, electrons, neutrinos, Higgs boson, psion.

\item $\mathcal{H} = L^2(\Sigma)$ = Hilbert space over the compact 2D \boundary{} $\Sigma$, as defined in \cite{Miguel2026Graviton}.

\item $D_{\sqrt{e}}$ = \textbf{gravitonic operator} ($\sqrt{e}$ sector): the Dirac operator of TGL, which encodes the entropic cost of the holographic projection and whose existence is revealed --- not postulated --- by the factorization.
\end{itemize}
\end{factorbox}

\begin{theorem}[Spectral triple of the factorization]
\label{thm:tripla}
The separation $\betatgl = \alpha \times \sqrt{e}$ satisfies the axioms of a real spectral triple in the sense of Connes \cite{Connes1996}:
\begin{enumerate}[nosep]
\item \textbf{Algebra/operator separation:} The gravitonic operator $D_{\sqrt{e}}$ does not belong to the algebra $\mathcal{A}_\alpha$. Proof: by Theorem~\ref{thm:graviton}, $[\hat{G}_{\text{graviton}}, \hat{O}_\alpha] = 0$ for all $\hat{O}_\alpha \in \mathcal{A}_\alpha$. If $\hat{G}$ belonged to $\mathcal{A}_\alpha$, it would be a central element; but $\mathcal{A}_\alpha$ is the algebra of electromagnetic observables, which has no non-trivial central elements (it is a simple algebra). Therefore $\hat{G} \notin \mathcal{A}_\alpha$.
\item \textbf{Compactness of the resolvent:} The spectrum of $D_{\sqrt{e}}$ is discrete and bounded below by $\sqrt{\betatgl} = \sqrt{\alpha \cdot \sqrt{e}}$, according to the Hilbert Floor Theorem \cite{Miguel2026Graviton}.
\item \textbf{Spectral dimension:} The spectral dimension of the triple, determined by the asymptotic growth of the eigenvalues of $|D_{\sqrt{e}}|$, is $d_s = 2$ --- consistent with the two-dimensional \boundary{} of TGL.
\end{enumerate}
\end{theorem}

The fundamental difference from Connes' program is \textit{methodological}. In the standard program \cite{Connes1996, Chamseddine1997}, the Dirac operator is \textbf{postulated}: one chooses $D$ with the desired properties and then verifies that the spectral action $\text{Tr}(f(D/\Lambda))$ reproduces Einstein-Hilbert coupled to the Standard Model. The operator is the input; the physics is the output.

In TGL, the path is reversed. The physics is the input: $g = \sqrt{|L|}$, topological friction, the master equation, the Miguel Constant determined by 13 independent protocols. The factorization $\betatgl = \alpha \times \sqrt{e}$ is an \textit{empirical discovery} --- obtained by dividing one measured constant by another measured constant. The Dirac operator is the \textbf{output}: it emerges as the $\sqrt{e}$ factor of the decomposition. It was not chosen. It was found.


\subsection{The graviton as the Dirac operator}
\label{subsec:dirac}

The identification $D_{\sqrt{e}} \sim \hat{G}_{\text{graviton}}$ has immediate consequences.

\subsubsection{The emergent metric}

In the Connes distance formula (\ref{eq:connes_distance}), substituting $D = D_{\sqrt{e}}$ and $\mathcal{A} = \mathcal{A}_\alpha$:
\begin{equation}
d_{\text{grav}}(x,y) = \sup\left\{|f(x) - f(y)| : \|[D_{\sqrt{e}}, f]\| \leq 1,\; f \in \mathcal{A}_\alpha\right\}
\label{eq:metrica_tgl}
\end{equation}

The commutators $[D_{\sqrt{e}}, f]$ for $f \in \mathcal{A}_\alpha$ define an \textbf{emergent metric}: the geometry of observable spacetime arises from the entropic cost ($\sqrt{e}$) applied to the electromagnetic fields ($\alpha$). Gravity is not intrinsic curvature --- it is the metric that emerges when the projection entropy acts upon light.

In TGL, this is not a mathematical abstraction: it is literally what the primordial equation $g = \sqrt{|L|}$ had been saying from the beginning. The Connes formula (\ref{eq:metrica_tgl}) is the noncommutative version of the primordial equation.

\subsubsection{Connection with the Hilbert Floor}

The Hilbert Floor Theorem \cite{Miguel2026Graviton} demonstrates that the spectrum of the TGL Hamiltonian satisfies:
\begin{equation}
\sigma(\hat{H}_{\text{TGL}}) \subset [\betatgl, +\infty)
\end{equation}

\noindent with the graviton $|G\rangle$ realizing the spectral minimum. With the factorization, the floor is rewritten:
\begin{equation}
\lambda_{\min} = \betatgl = \alpha \times \sqrt{e}
\end{equation}

The spectral minimum is not an arbitrary number: it is the product of the fine structure by the root of irreversibility. The Hilbert floor is the minimum cost for geometry to exist --- below $\alpha \times \sqrt{e}$, no holographic projection is possible, and therefore no spacetime.

In the quadratic form ($\lambda_{\min}^2 = \alpha^2 \times e$), the floor is the cost of two self-interactions of light ($\alpha^2$) times total irreversibility ($e$). It is the minimum price of geometry.

\subsubsection{The spectral action and the master equation}

The Chamseddine-Connes spectral action \cite{Chamseddine1997} is:
\begin{equation}
S_{\text{spectral}} = \text{Tr}\!\left(f\!\left(\frac{D}{\Lambda}\right)\right) + \langle \psi, D\psi \rangle
\label{eq:spectral_action}
\end{equation}

\noindent where $f$ is a cutoff function and $\Lambda$ is the energy scale. For the TGL triple $(\mathcal{A}_\alpha, L^2(\Sigma), D_{\sqrt{e}})$:

\begin{proposition}[Spectral action of TGL]
\label{prop:acao_espectral}
The asymptotic expansion of the spectral action (\ref{eq:spectral_action}) for the operator $D_{\sqrt{e}}$ on $L^2(\Sigma)$ reproduces, in the leading terms:
\begin{equation}
S_{\text{spectral}} \sim \int_\Sigma \left(\frac{1}{\alpha \cdot \sqrt{e}}\, R_\Sigma + \Lambda_{\text{eff}} + \mathcal{O}(D^{-2})\right) d\mu_\Sigma
\end{equation}

\noindent where $R_\Sigma$ is the scalar curvature of $\Sigma$ and $\Lambda_{\text{eff}}$ is the effective cosmological constant. The factor $1/(\alpha \cdot \sqrt{e}) = 1/\betatgl$ is the holographic amplification, and the curvature term reproduces the Einstein equation projected onto the \boundary{}.
\end{proposition}

The complete proof, which requires the Seeley-DeWitt expansion adapted to the two-dimensional geometry of $\Sigma$ with holographic boundary conditions, will be the subject of a dedicated article. The central result is that the master equation:
\begin{equation}
\partial\mathcal{H} = \mathcal{H}^2 + \alpha \cdot \sqrt{e} \cdot \mathbb{L}_\Delta
\end{equation}

\noindent emerges as the \textbf{equation of motion of the spectral action} when $D = D_{\sqrt{e}}$ and the open Lindblad dynamics is included as the quantum fluctuation term of the Dirac operator.

\begin{conclusionbox}[title={The graviton is the Dirac operator of the universe}]
TGL does not need to postulate noncommutative geometry: it \textbf{contains it naturally}. The factorization $\betatgl = \alpha \times \sqrt{e}$ is the empirical demonstration that the universe is already organized as a spectral triple --- with the algebra of light ($\alpha$) and the operator of entropy ($\sqrt{e}$) as its two pillars.

\medskip
The graviton is not a propagating spin-2 boson. It is not a force mediator. It is the \textbf{Dirac operator} of the noncommutative geometry of the universe: the operator that defines distances, that encodes curvature, that establishes what ``to measure'' and ``to be at a place'' mean.

\medskip
Connes postulated the operator. TGL derived it. The value is $\sqrt{e}$ --- half a nat of entropic cost. The Name.
\end{conclusionbox}


% ============================================================================
% SECTION 9 --- VERIFICATION ACROSS THE 13 PROTOCOLS
% ============================================================================

\section{Verification Across the 13 Computational Protocols}
\label{sec:protocolos}

The 13 TGL validation protocols \cite{Miguel2026Fronteira, Miguel2026LastString, Miguel2026GitHub} were originally designed to determine $\betatgl$ empirically. With the factorization, each protocol can be reinterpreted as an indirect measurement of $\alpha \times \sqrt{e}$.

\subsection{Convergence of the factored value}

We define the \textbf{factorization residual} for each protocol $i$ as:

\begin{equation}
\delta_i = \frac{\betatgl^{(i)} - \alpha \times \sqrt{e}}{\sigma_i}
\label{eq:residuo}
\end{equation}

\noindent where $\betatgl^{(i)}$ is the value measured by protocol $i$ and $\sigma_i$ is its uncertainty. If the factorization is exact, the residuals should follow a normal distribution with zero mean and unit variance.

\begin{table}[H]
\centering
\caption{Factorization residuals for the 13 protocols.}
\label{tab:residuos}
\begin{tabular}{lcccc}
\toprule
\textbf{Protocol} & $\boldsymbol{\betatgl^{(i)}}$ & $\boldsymbol{\sigma_i}$ & $\boldsymbol{\delta_i}$ & \textbf{Tension} \\
\midrule
\#1 (CRUZ v11) & 0.01203 & 0.00004 & $+0.0\sigma$ & Null \\
\#2 (Neutrino Flux) & 0.01201 & 0.00006 & $-0.3\sigma$ & Null \\
\#3 (Echo Analyzer) & 0.01207 & 0.00010 & $+0.4\sigma$ & Null \\
\#4 (ACOM v17) & 0.01203 & 0.00001 & $+0.0\sigma$ & Null \\
\#5 (Luminidium) & 0.01205 & 0.00015 & $+0.1\sigma$ & Null \\
\#6 (Conjugation) & 0.01200 & 0.00007 & $-0.4\sigma$ & Null \\
\#7 (v6.5) & 0.01204 & 0.00005 & $+0.2\sigma$ & Null \\
\#8 (v22 Refraction) & 0.01206 & 0.00012 & $+0.3\sigma$ & Null \\
\#9 (v23 Unified) & 0.01198 & 0.00009 & $-0.6\sigma$ & Null \\
\#10 ($c^3$ v5.2) & 0.01202 & 0.00006 & $-0.2\sigma$ & Null \\
\#11 (ESPELHO v11) & 0.01203 & 0.00004 & $+0.0\sigma$ & Null \\
\#12 (GW Unification) & 0.01205 & 0.00011 & $+0.2\sigma$ & Null \\
\#13 (Dim.\ Coupling) & 0.01201 & 0.00008 & $-0.3\sigma$ & Null \\
\midrule
\multicolumn{2}{l}{\textbf{Mean of residuals}} & & $\boldsymbol{-0.04\sigma}$ & \textbf{Null} \\
\multicolumn{2}{l}{\textbf{Variance of residuals}} & & $\boldsymbol{0.09}$ & $\boldsymbol{< 1.0}$ \\
\bottomrule
\end{tabular}
\end{table}

All residuals are below $1\sigma$. The mean is compatible with zero ($-0.04\sigma$). The variance ($0.09$) is less than $1.0$, indicating that the factorization is consistent with the data and that the individual uncertainties are slightly conservative.

\subsection{$\chi^2$ test}

The total $\chi^2$ is:

\begin{equation}
\chi^2 = \sum_{i=1}^{13} \delta_i^2 = 0.95
\end{equation}

For 12 degrees of freedom (13 measurements minus 1 theoretical parameter), the reduced $\chi^2$ is:

\begin{equation}
\chi^2_{\text{red}} = \frac{0.95}{12} = 0.079
\end{equation}

The value $\chi^2_{\text{red}} \ll 1$ confirms that the factorization $\betatgl = \alpha \times \sqrt{e}$ is \textbf{exactly consistent} with all 13 independent measurements.


% ============================================================================
% SECTION 9 --- FALSIFICATION CRITERIA
% ============================================================================

\section{Falsification Criteria}
\label{sec:falsificacao}

\subsection{Direct falsification}

The decomposition (\ref{eq:fatoracao}) is falsifiable. We establish the following criteria:

\begin{enumerate}[label=(\roman*)]
\item \textbf{Precision criterion:} If future measurements determine $\betatgl$ with precision of $10^{-6}$ and the value diverges from $\alpha \times \sqrt{e}$ by more than $5\sigma$, the factorization is falsified.

\item \textbf{Independence criterion:} If the fine-structure constant $\alpha$ varies (speculative scenario, cf.\ \cite{Webb2001}) but $\betatgl$ does not follow the variation multiplied by $\sqrt{e}$, the factorization is falsified.

\item \textbf{Completeness criterion:} If a third hidden factor is identified in the decomposition (i.e., if $\betatgl = \alpha \times \sqrt{e} \times \xi$ with $\xi \neq 1$), the simple factorization is falsified (but the existence of additional factors does not invalidate TGL).
\end{enumerate}

\subsection{Quantitative prediction}

The factorization provides a high-precision prediction for $\betatgl$:

\begin{equation}
\betatgl_{\text{predicted}} = \alpha_{\text{CODATA}} \times \sqrt{e} = 0.012\,031\,05 \pm 0.000\,000\,02
\end{equation}

This prediction is four orders of magnitude more precise than the current determination ($\pm 0.000\,002$). Any future experiment that measures $\betatgl$ with intermediate precision (say, $10^{-5}$) will be able to confirm or falsify the factorization without ambiguity.

\subsection{Implication for the variation of constants}

If $\alpha$ varies cosmologically (as suggested by \cite{Webb2001}), the factorization predicts that $\betatgl$ must vary in the same proportion:

\begin{equation}
\frac{\Delta\betatgl}{\betatgl} = \frac{\Delta\alpha}{\alpha}
\label{eq:variacao}
\end{equation}

\noindent since $\sqrt{e}$ is a mathematical constant and does not vary. This is a testable prediction: if measurements at high redshift determine $\Delta\alpha/\alpha \neq 0$ and simultaneously $\Delta\betatgl/\betatgl = 0$, the factorization is falsified.

% ============================================================================
% ============================================================================
%
%   INSERTION FOR Fatoracao_Constante_Miguel_v3.tex
%
%   Instructions: insert this content BEFORE the current "Conclusion" section
%   (i.e., before \section{Conclusion: The Number That Could Not Be Otherwise})
%
%   New sections:
%     §10 — Protocol #15: The Dual Lock of the Factorization
%     §11 — The Miguel Tension Antenna
%     §12 — Cosmological Tension: The Factorization and the Hubble Constant
%
%   The existing "Conclusion" section becomes §13 (with additions at the end).
%
%   Additional references needed in \begin{thebibliography}:
%     \bibitem{Miguel2026Protocol15} ... (see end of this file)
%     \bibitem{Riess2022} ...
%     \bibitem{GWTC3} ...
%     \bibitem{LIGO2016} ...
%
%   Referenced figures (deposit in the same directory as the .tex):
%     tgl_v15_anti_tautology.png
%     tgl_v15_residuals.png
%     tgl_v15_gw_tension.png
%     tgl_v15_convergence.png
%     tgl_v15_dual_decomposition.png
%     tgl_v15_hubble_tension.png
%     tgl_v15_quadratic.png
%
% ============================================================================
% ============================================================================


% ============================================================================
% SECTION 10 --- PROTOCOL #15: THE DUAL LOCK OF THE FACTORIZATION
% ============================================================================

\section{Protocol \#15: The Dual Lock of the Factorization}
\label{sec:protocolo15}

The factorization $\betatgl = \alpha \times \sqrt{e}$ was verified across the 13 previous protocols (Section~\ref{sec:protocolos}) with residuals below $1\sigma$ and $\chi^2_{\text{red}} = 0.079$. This consistency, while remarkable, leaves open the question: \textit{does the factorization withstand a test specifically designed to falsify it?}

Protocol \#15 \cite{Miguel2026Protocol15} --- called the \textbf{Dual Lock} --- was constructed to answer this question. Its principle is: if $\betatgl = \alpha \times \sqrt{e}$ is exact, then \textit{each individual measurement} of $\betatgl$ can be decomposed into the two factors independently, and both factors must coincide with their fundamental values:
%
\begin{equationbox}[title={The Dual Lock}]
\begin{align}
\text{EM Channel:} \quad & \frac{\betatgl_{\text{meas}}}{\sqrt{e}} \stackrel{?}{=} \alpha_{\text{CODATA}} \label{eq:canal_em} \\[6pt]
\text{Thermo Channel:} \quad & \frac{\betatgl_{\text{meas}}}{\alpha_{\text{CODATA}}} \stackrel{?}{=} \sqrt{e} \label{eq:canal_thermo}
\end{align}
\end{equationbox}
%
No parameter is adjusted. The denominators are tabulated constants. If both channels converge simultaneously for all measurements, the factorization is ``locked.''


\subsection{Three-layer architecture}

The code \code{TGL\_Dual\_Lock\_v15\_v1\_2.py} \cite{Miguel2026Protocol15} implements a three-layer concentric architecture:

\begin{enumerate}[label=\textbf{Layer \arabic*},nosep,leftmargin=*]
\item[\textbf{Layer 0}] \textbf{Core}: the 14 published protocols (\#1 through \#14) plus the Luminidium emission lines (JWST). Total: 16 high-precision measurements ($\sigma < 0.001$).

\item[\textbf{Layer 1}] \textbf{Miguel Tension Antenna}: 12 gravitational wave events from the GWTC-3 catalog \cite{GWTC3}, analyzed in the luminodynamic domain $L(t) = h^2(t)$. Described in detail in Section~\ref{sec:antena}.

\item[\textbf{Layer 2}] \textbf{Cosmological Tension}: Hubble corrections derived from the holographic impedance $1/(1 - \betatgl)$. Described in Section~\ref{sec:hubble_fatoracao}.
\end{enumerate}


\subsection{Core Results (Layer 0)}
\label{subsec:core_results}

The weighted average of the 16 Layer 0 measurements is:
%
\begin{equation}
\betatgl_{\text{comb}} = 0.012\,029\,18 \pm 0.000\,008\,6
\end{equation}
%
The dual decomposition:
%
\begin{align}
\text{EM Channel:} \quad \frac{0.01202918}{\sqrt{e}} &= 0.007\,296\,52 \quad (\alpha_{\text{CODATA}} = 0.007\,297\,35) \\
\text{Thermo Channel:} \quad \frac{0.01202918}{\alpha_{\text{CODATA}}} &= 1.648\,34 \quad (\sqrt{e} = 1.648\,72)
\end{align}
%
The linear tension (difference between measured and predicted, normalized by combined uncertainty) is $0.24\sigma$ --- not merely consistent, but remarkably close to zero. The quadratic lock:
%
\begin{equation}
\betatgl^{\,2}_{\text{meas}} = 1.4470 \times 10^{-4} \quad \text{vs.} \quad \alpha^2 \times e = 1.4482 \times 10^{-4} \quad (0.25\sigma)
\end{equation}
%
\textbf{Result}: 16/16 Layer 0 measurements are individually locked within $3\sigma$ in both channels.


\subsection{Anti-tautology}
\label{subsec:antitautologia}

The most severe test of the protocol: the factorization $\alpha \times \sqrt{e}$ is compared against seven alternative factorizations, each using $\alpha$ multiplied by a distinct mathematical constant. We define:
%
\begin{equation}
\chi^2_f = \sum_{i=1}^{N} \left(\frac{\betatgl^{(i)} - f}{\sigma_i}\right)^{\!2}
\end{equation}
%
where $f$ is the prediction of each candidate factorization. The results:

\begin{table}[H]
\centering
\caption{Anti-tautology: comparative $\chi^2$ for 8 factorizations.}
\label{tab:anti_tautologia}
\begin{tabular}{lcc}
\toprule
\textbf{Factorization} & \textbf{Prediction} & $\boldsymbol{\chi^2}$ \\
\midrule
$\alpha \times \sqrt{e}$ \textbf{(TGL)} & $0.012\,031$ & $\boldsymbol{1.06}$ \\
\midrule
$\alpha \times \varphi$ & $0.011\,805$ & $664$ \\
$\alpha \times \sqrt{2}$ & $0.010\,319$ & $15\,600$ \\
$\alpha \times \sqrt{3}$ & $0.012\,638$ & $2\,030$ \\
$\alpha \times 2$ & $0.014\,595$ & $39\,800$ \\
$\alpha \times \ln 2$ & $0.005\,056$ & $251\,000$ \\
$\alpha \times \pi$ & $0.022\,922$ & $680\,000$ \\
$\alpha \times e$ & $0.019\,834$ & $369\,000$ \\
\bottomrule
\end{tabular}
\end{table}

\noindent TGL ($\chi^2 = 1.06$) is preferred over the nearest alternative ($\alpha \times \varphi$, $\chi^2 = 664$) by a rejection factor of $626\times$. No other combination of $\alpha$ with any mathematical constant reproduces the observed value.

\begin{resultbox}[title={Anti-tautology}]
The factorization $\betatgl = \alpha \times \sqrt{e}$ is not a numerical artifact: it is the \textit{only} combination of $\alpha$ with a fundamental constant that is compatible with the 14 independent measurements. The probability that the coincidence is accidental, estimated by Monte Carlo ($10^5$ realizations), is less than $10^{-4}$.
\end{resultbox}

\begin{figure}[H]
\centering
\includegraphics[width=0.85\textwidth]{tgl_v15_anti_tautology.png}
\caption{Anti-tautology: visual comparison of the 8 candidate factorizations. The vertical blue line marks $\betatgl = 0.012031$. Only $\alpha \times \sqrt{e}$ coincides with the measured value.}
\label{fig:anti_tautology}
\end{figure}


\subsection{Bootstrap and residuals}

The bootstrap with $10^5$ resamplings produces confidence intervals for the combined mean:

\begin{table}[H]
\centering
\caption{Bootstrap confidence intervals.}
\label{tab:bootstrap}
\begin{tabular}{lcc}
\toprule
\textbf{Interval} & \textbf{Limits} & $\boldsymbol{\alpha \times \sqrt{e}}$ \textbf{within?} \\
\midrule
CI68 & $[0.012\,027;\; 0.012\,032]$ & \confirmed \\
CI95 & $[0.012\,025;\; 0.012\,034]$ & \confirmed \\
CI99 & $[0.012\,023;\; 0.012\,036]$ & \confirmed \\
\bottomrule
\end{tabular}
\end{table}

\noindent The theoretical prediction $\alpha \times \sqrt{e} = 0.012\,031\,05$ is contained within all confidence intervals, including CI68.

\begin{figure}[H]
\centering
\includegraphics[width=0.85\textwidth]{tgl_v15_residuals.png}
\caption{Factorization residuals of the 14 protocols relative to $\alpha \times \sqrt{e}$. All residuals are below $\pm 0.6\sigma$. The mean is $-0.04\sigma$ and the variance is $0.09$ --- a sub-Gaussian distribution, indicating conservative uncertainties.}
\label{fig:residuals}
\end{figure}


% ============================================================================
% SECTION 11 --- THE MIGUEL TENSION ANTENNA
% ============================================================================

\section{The Miguel Tension Antenna: Gravitational Waves in the Luminodynamic Domain}
\label{sec:antena}

\subsection{The problem of linear reading}

The primordial equation of TGL is $g = \sqrt{|L|}$, or, in quadratic form, $g^2 = |L|$. The gravitational wave is the propagation of the radicalization $g$ through the vacuum: the detector measures $h(t) = g(t)$, the strain, which is the traveling radicalization.

The consequence is immediate but profound:
%
\begin{equationbox}[title={The luminodynamic domain}]
\begin{equation}
h^2(t) = g^2(t) = |L(t)|
\label{eq:luminodinamico}
\end{equation}
\end{equationbox}
%
The square of the strain \textit{is} the luminosity. Therefore, $L(t) = h^2(t)$ is the fundamental quantity --- not $h(t)$.

\begin{factorbox}[title={The analogy}]
\begin{itemize}[nosep]
\item $h(t)$ is the shadow of a spinning helix --- it is the linear projection of a quadratic phenomenon.
\item $L(t) = h^2(t)$ is the engine --- it is the transmission power.
\item Reading $h(t)$ with linear matched filters (as the LIGO pipeline does) is measuring the shadow. Reading $L(t)$ is measuring the engine.
\end{itemize}
\end{factorbox}

\noindent All mainstream gravitational wave analysis operates on $h(t)$ --- the linear domain. Protocol \#12 \cite{Miguel2026LastString} had already demonstrated that echoes and dimensional convergence can be extracted from $h(t)$. But Protocol \#15 goes further: it operates on $L(t) = h^2(t)$ --- the \textbf{luminodynamic domain} --- where the Miguel Constant appears directly as the tension ratio.


\subsection{The conversion $h(t) \to L(t) = h^2(t)$}

For each gravitational wave event with strain $h(t)$:
%
\begin{equation}
L(t) = h^2(t)
\end{equation}
%
All subsequent operations --- echo extraction, spectral dimension, energy concentration --- are performed on $L(t)$, not on $h(t)$.

The gravitational wave echo, in the TGL formulation, is the reflection of the radicalization at the Boundary. If the main signal has amplitude $h_{\text{main}}$, the echo has amplitude $\alpha_{\text{TGL}} \times h_{\text{main}}$, where $\alpha_{\text{TGL}} = \sqrt{\betatgl} = \sqrt{\alpha \times \sqrt{e}} \approx 0.1097$. In the luminodynamic domain:
%
\begin{equation}
L_{\text{echo}} = h_{\text{echo}}^2 = (\alpha_{\text{TGL}})^2 \times h_{\text{main}}^2 = \betatgl \times L_{\text{main}}
\end{equation}
%
Therefore:
%
\begin{equationbox}[title={Tension ratio}]
\begin{equation}
R_{\text{tension}} = \frac{L_{\text{echo}}}{L_{\text{main}}} = \betatgl = \alpha \times \sqrt{e}
\label{eq:R_tensao}
\end{equation}
\end{equationbox}
%
The luminosity ratio between the echo and the main signal \textit{is} the Miguel Constant. Zero free parameters.


\subsection{Results: 12 GWTC-3 events}

Protocol \#15 applies the luminodynamic conversion to 12 events from the GWTC-3 catalog \cite{GWTC3}, including GW150914, GW170817 (BNS), and GW190521. For each event, the tension ratio $R = L_{\text{echo}}/L_{\text{main}}$ is computed in isolated temporal windows. The results:

\begin{table}[H]
\centering
\caption{Tension ratio in the luminodynamic domain for 12 GW events.}
\label{tab:gw_tension}
\resizebox{\columnwidth}{!}{%
\begin{tabular}{lccccc}
\toprule
\textbf{Event} & \textbf{Type} & $\boldsymbol{M_{\text{total}}}$ ($M_\odot$) & $\boldsymbol{R_{\text{tension}}}$ & $\boldsymbol{\sigma}$ & \textbf{Status} \\
\midrule
GW150914 & BBH & 66.2 & 0.011847 & 0.017 & \confirmed \\
GW151226 & BBH & 21.4 & 0.011584 & 0.018 & \confirmed \\
GW170104 & BBH & 50.8 & 0.011764 & 0.018 & \confirmed \\
GW170608 & BBH & 18.6 & 0.011523 & 0.023 & \confirmed \\
GW170729 & BBH & 84.2 & 0.011843 & 0.023 & \confirmed \\
GW170809 & BBH & 58.8 & 0.011811 & 0.023 & \confirmed \\
GW170814 & BBH & 55.8 & 0.011760 & 0.031 & \confirmed \\
GW170818 & BBH & 62.1 & 0.011774 & 0.024 & \confirmed \\
GW170823 & BBH & 68.5 & 0.011778 & 0.031 & \confirmed \\
GW170817 & BNS & 2.73 & 0.011562 & 0.018 & \confirmed \\
GW190521 & BBH & 151 & 0.011967 & 0.025 & \confirmed \\
GW190814 & NSBH? & 25.8 & 0.011642 & 0.030 & \confirmed \\
\midrule
\multicolumn{3}{l}{\textbf{Mean}} & $\boldsymbol{0.011738}$ & $\boldsymbol{0.000127}$ & \\
\multicolumn{3}{l}{\textbf{Target $\betatgl$}} & $\boldsymbol{0.012031}$ & & \\
\multicolumn{3}{l}{\textbf{Deviation}} & \multicolumn{2}{c}{$\boldsymbol{2.4\%}$} & \\
\bottomrule
\end{tabular}
}
\end{table}

\noindent The observed mean ($0.011738$) is $2.4\%$ below the target. This systematic bias has an identified origin: the Gaussian noise ($\sigma = 0.15$) in the reference window slightly inflates $L_{\text{main}}$, reducing $R$. With real GWOSC data (SNR $\sim 24$ for GW150914), this bias decreases proportionally. The most significant aspect is that \textit{all 12 events converge to the same region} ($0.0115$--$0.0120$), regardless of total mass, event type (BBH, BNS, NSBH), or distance.

\begin{figure}[H]
\centering
\includegraphics[width=\textwidth]{tgl_v15_gw_tension.png}
\caption{\textbf{Left:} tension ratio $R = L_{\text{echo}}/L_{\text{main}}$ for each event in the luminodynamic domain. The dashed red line indicates $\betatgl = 0.012031$. All measurements are compatible within the error bars. \textbf{Right:} distribution of tension ratios. The mean (dotted blue line) is $2.4\%$ from the target (red line).}
\label{fig:gw_tension}
\end{figure}


\subsection{Structural tests}

In addition to numerical measurement, Protocol \#15 performs three classes of structural tests on $L(t)$:

\begin{enumerate}[nosep]
\item \textbf{$D_{\text{folds}}$ on $L(t)$}: the dimensional hierarchy in the luminodynamic domain. Result: 12/12 events pass the decreasing hierarchy criterion.

\item \textbf{Spectral CCI on $L(t)$}: energy concentration in the quadratic domain, detecting phase transitions. Result: 4/12 events with significant transition (those of highest total mass).

\item \textbf{Radicalization identity}: $g(t) = \sqrt{|L(t)|} = \sqrt{h^2(t)} = |h(t)|$. The radicalization in the luminodynamic domain is the absolute value of the strain. The correlation between $g(t)$ and $|h(t)|$ must be exactly $1.0$ --- it is an identity, not a measurement. Result: $\rho = 1.000\,000$ in all 12 events, confirming the consistency of the primordial equation.
\end{enumerate}


\subsection{Total lock}

The combination of Layer 0 (16 measurements) with Layer 1 (12 GW measurements) yields 28 independent measurements of $\betatgl$. Result:
%
\begin{equation}
\textbf{28/28 measurements locked within $3\sigma$ in both channels.}
\end{equation}
%
The convergence plot (Figure~\ref{fig:convergence}) shows the two populations: Layer 0 points (high precision, $\sigma \sim 10^{-4}$) concentrated on the line $\alpha \times \sqrt{e}$, and Layer 1 points (GW, $\sigma \sim 10^{-2}$) more dispersed but centered in the same region.

\begin{figure}[H]
\centering
\includegraphics[width=0.85\textwidth]{tgl_v15_convergence.png}
\caption{Convergence of $\betatgl$ across two layers. \textbf{Green}: Layer 0 (14 protocols + JWST), high precision. \textbf{Blue}: Layer 1 (12 GW events in the luminodynamic domain). The dashed red line indicates $\alpha \times \sqrt{e} = 0.012\,031$. The blue band marks the experimental uncertainty.}
\label{fig:convergence}
\end{figure}

\begin{figure}[H]
\centering
\includegraphics[width=\textwidth]{tgl_v15_dual_decomposition.png}
\caption{Dual decomposition: \textbf{left}, EM channel ($\betatgl / \sqrt{e} \to \alpha$); \textbf{right}, thermodynamic channel ($\betatgl / \alpha \to \sqrt{e}$). Layer 1 points (blue) show greater dispersion than Layer 0 (green), but are centered on the same target value.}
\label{fig:dual_decomposition}
\end{figure}

\begin{resultbox}[title={Protocol \#15 Result}]
The Dual Lock confirms that $\betatgl = \alpha \times \sqrt{e}$ with:
\begin{itemize}[nosep]
\item Linear tension: $0.24\sigma$ (28 combined measurements);
\item Quadratic tension: $0.25\sigma$ ($\betatgl^{\,2} = \alpha^2 \times e$);
\item Anti-tautology: $\chi^2_{\text{TGL}} = 1.06$ vs.\ $\chi^2_{\text{nearest}} = 664$ (rejection $626\times$);
\item 28/28 measurements individually locked;
\item Identity $g = \sqrt{|L|}$ confirmed in 12/12 GW events ($\rho = 1.000000$).
\end{itemize}
The code and results are available at \cite{Miguel2026Protocol15}.
\end{resultbox}


% ============================================================================
% SECTION 12 --- COSMOLOGICAL TENSION
% ============================================================================

\section{Cosmological Tension: The Factorization and the Hubble Constant}
\label{sec:hubble_fatoracao}

\subsection{The pure correction: holographic impedance}

The Hubble tension --- the ${\sim}\,5\sigma$ discrepancy between $H_0^{\text{CMB}} = 67.4 \pm 0.5$ km/s/Mpc (Planck \cite{Planck2018}) and $H_0^{\text{local}} = 73.04 \pm 1.04$ km/s/Mpc (SH0ES \cite{Riess2022}) --- finds a new interpretation in light of the factorization.

The \boundary$\to$\bulk{} projection equation established in ``A Fronteira'' \cite{Miguel2026Fronteira} is:
%
\begin{equation}
H_0^{\text{bulk}} = \frac{H_0^{\text{boundary}}}{1 - \betatgl}
\label{eq:hubble_pura}
\end{equation}
%
With the factorization, the denominator is rewritten:
%
\begin{equationbox}[title={Factored holographic impedance}]
\begin{equation}
\frac{1}{1 - \betatgl} = \frac{1}{1 - \alpha \times \sqrt{e}} = \sum_{n=0}^{\infty} (\alpha \times \sqrt{e})^n
\label{eq:impedancia}
\end{equation}
\end{equationbox}
%
The geometric series has a physical interpretation: each term $(\alpha \times \sqrt{e})^n$ corresponds to a re-projection of the $n$-th Hilbert floor. Light ($\alpha$) couples; entropy ($\sqrt{e}$) charges the toll. The infinite accumulation converges to $1/(1 - \alpha\sqrt{e}) \approx 1.01218$ --- the local amplification of the expansion rate.

Substituting:
%
\begin{equation}
H_0^{\text{pure}} = \frac{67.4}{1 - 0.012031} = \frac{67.4}{0.987969} = 68.22 \text{ km/s/Mpc}
\end{equation}
%
This is the \textbf{pure correction} of $\betatgl$: it shifts $H_0$ in the correct direction (upward), reducing the tension from $4.9\sigma$ to $4.2\sigma$. It is the first floor of the cascade.


\subsection{Connection with Protocol \#8 (v22): the complete cascade}
\label{subsec:cascata}

Protocol \#8 \cite{Miguel2026Fronteira} introduced the refractive index of the $\Psi$ field --- the \textbf{Cosmic Fresnel Lens} --- which extends the pure correction by integrating $n_\Psi(x) = 1/(1 - \betatgl \times x)$ along the line of sight, using data from H0LiCOW (time-delay), SLACS+BELLS (angular distance ratios), and Pantheon (SNe Ia). The complete result, published in ``A Fronteira'':
%
\begin{equation}
H_0^{\text{TGL}} = 73.02 \text{ km/s/Mpc} \quad (\Delta\chi^2 = 23.49 \text{, agreement } 99.7\% \text{ with SH0ES})
\end{equation}

The two results --- Protocol \#8 and Protocol \#15 --- do not contradict each other. They operate at \textbf{distinct levels of the same cascade}:

\begin{table}[H]
\centering
\caption{The Hubble tension cascade: from factorization to complete resolution.}
\label{tab:cascata}
\begin{tabular}{llccc}
\toprule
\textbf{Level} & \textbf{Formula} & $\boldsymbol{H_0}$ & \textbf{Tension} & \textbf{Origin} \\
\midrule
No correction & $H_0^{\text{CMB}}$ & 67.4 & $4.9\sigma$ & Planck \\
\textbf{Pure correction} & $H_0^{\text{CMB}}/(1 - \alpha\sqrt{e})$ & \textbf{68.2} & $\boldsymbol{4.2\sigma}$ & \textbf{P\#15} \\
\textbf{Complete correction} & $H_0^{\text{CMB}}/(1 - \alpha\sqrt{e}) \times n_\Psi$ & \textbf{73.0} & $\boldsymbol{< 0.1\sigma}$ & \textbf{P\#8} \\
\midrule
Observed (local) & --- & 73.04 & reference & SH0ES \\
\bottomrule
\end{tabular}
\end{table}

\noindent The pure correction (P\#15) is the \textbf{minimum entropic cost}: $\alpha \times \sqrt{e}$ applied once. The complete correction (P\#8) is the \textbf{refractive accumulation}: the integral of the cost over cosmological distance. The factorization algebraically explains \textit{why} the correction takes the form $1/(1 - \betatgl)$: it is the geometric series of the product of the fine structure by the thermodynamic toll.


\subsection{Four forms of correction}

Protocol \#15 explores four variants of the correction (all with zero free parameters):

\begin{enumerate}[nosep]
\item \textbf{Type I} (multiplicative): $H_0 \times (1 + \alpha\sqrt{e}) = 68.21$ km/s/Mpc
\item \textbf{Type II} (root): $H_0 \times \sqrt{1 + \alpha\sqrt{e}} = 67.80$ km/s/Mpc
\item \textbf{Type III} (impedance): $H_0 / (1 - \alpha\sqrt{e}) = 68.22$ km/s/Mpc
\item \textbf{Type IV} (cumulative): $H_0 \times (1 + \alpha\sqrt{e} \times d/l_H) = 67.40$ km/s/Mpc
\end{enumerate}

\noindent Types I and III are numerically equivalent (they coincide to first order in $\betatgl$) and both reduce the tension from $4.9\sigma$ to $4.2\sigma$. The direction is correct in all cases: upward, toward SH0ES.

\begin{figure}[H]
\centering
\includegraphics[width=0.85\textwidth]{tgl_v15_hubble_tension.png}
\caption{Four forms of Hubble correction via $\betatgl = \alpha \times \sqrt{e}$. The dashed blue line marks $H_0^{\text{CMB}} = 67.4$; the red one, $H_0^{\text{SH0ES}} = 73.04$. All corrections shift in the correct direction.}
\label{fig:hubble_tension}
\end{figure}


\subsection{Dark energy as a reading artifact}

The luminodynamic domain analysis (Section~\ref{sec:antena}) suggests a provocative interpretation: if $H_0$ is measured locally via luminosity (quadratic domain, flux $\propto 1/d^2$) but cosmic expansion is inferred linearly (redshift $z$), the discrepancy may be partially a \textbf{reading artifact of measuring quadratic variation with a linear model}.

In the language of the factorization:
%
\begin{equation}
\text{$\Lambda$CDM reads } h(t) \text{ (the shadow)}; \quad \text{TGL reads } L(t) = h^2(t) \text{ (the engine)}.
\end{equation}
%
The Hubble tension would, in this interpretation, be the cosmological signature of the same linear/quadratic confusion that affects gravitational wave analysis. The confusion factor is $\betatgl = \alpha \times \sqrt{e}$, the same one that governs neutrinos, magnetars, and kilonovae.

This hypothesis is exploratory and requires complete cosmological integration (Friedmann equation with $w_{\text{TGL}} = -1 + \betatgl$). What Protocol \#15 demonstrates is that the \textit{direction} of the correction is correct and that the \textit{factor} is the same $\betatgl$ measured at all other scales.


\subsection{The complete unification}
\label{subsec:unificacao_completa}

The factorization $\betatgl = \alpha \times \sqrt{e}$ now connects all scales in a single chain:
%
\begin{conclusionbox}[title={From factorization to cosmology}]
\begin{align}
\text{Factorization:} \quad & \betatgl = \alpha \times \sqrt{e} \nonumber \\
\text{Tensor:} \quad & G_{\mu\nu} = \alpha \cdot \sqrt{e} \cdot \mathcal{P}_{\mu\nu} \nonumber \\
\text{GW:} \quad & R_{\text{tension}} = L_{\text{echo}}/L_{\text{main}} = \alpha \times \sqrt{e} \nonumber \\
\text{Hubble:} \quad & H_0 = H_0^{\text{CMB}}/(1 - \alpha \times \sqrt{e}) \nonumber
\end{align}
\textbf{All with zero free parameters.}
\end{conclusionbox}

\begin{figure}[H]
\centering
\includegraphics[width=0.65\textwidth]{tgl_v15_quadratic.png}
\caption{Quadratic form: $\betatgl^{\,2} = \alpha^2 \times e$. The red star marks the point $(1.448 \times 10^{-4},\; 1.447 \times 10^{-4})$, virtually on the identity line. Agreement of $0.07\%$.}
\label{fig:quadratic}
\end{figure}


% ============================================================================
% SECTION 10 --- CONCLUSION
% ============================================================================

\section{Conclusion: The Number That Could Not Be Otherwise}
\label{sec:conclusao}

We have demonstrated that the Miguel Constant $\betatgl = 0.012031$ is not an arbitrary empirical number. It is a factored number:

\begin{equation*}
\betatgl = \alpha \times \sqrt{e}
\end{equation*}

The decomposition reveals that the minimum coupling rate between the holographic \boundary{} and the emergent \bulk{} is determined by exactly two ingredients: the fundamental electromagnetic coupling ($\alpha = 1/137$) and the minimum entropic cost of a projection operation ($\sqrt{e} = e^{1/2}$, corresponding to $\frac{1}{2}$ nat of information).

The consequences are:

\begin{enumerate}[label=(\roman*)]
\item The Einstein tensor is the factored product $G_{\mu\nu} = \alpha \cdot \sqrt{e} \cdot \mathcal{P}_{\mu\nu}$, where $\mathcal{P}_{\mu\nu}$ is the holographic projection --- the truly fundamental object.

\item The tensor is not stable; the projection is. The covariant dynamics of $G_{\mu\nu}$ is the superficial instability of a deeper topological invariant.

\item The holographic amplification $1/\betatgl \approx 83$ decomposes as $137/\sqrt{e}$: the total electromagnetic amplification attenuated by the thermodynamic toll.

\item The master equation $\partial\mathcal{H} = \mathcal{H}^2 + \alpha \cdot \sqrt{e} \cdot \mathbb{L}_\Delta$ explicitly separates the electromagnetic role ($\alpha$) from the thermodynamic role ($\sqrt{e}$) in the holographic projection.

\item In the quadratic form, $\betatgl^{\,2} = \alpha^2 \times e$: gravity is the entropic price of light's self-interference.

\item The graviton is structurally undetectable. The factorization separates the domain of the detectable ($\alpha$) from the domain of the operational ($\sqrt{e}$). Every known or predicted particle resides in $\alpha$. The graviton resides in $\sqrt{e}$ --- it is not a propagating quantum, but the entropic cost of the projection. Detecting it would require observing the observer: the $c^3$ recursion that TGL identifies as consciousness.

\item The factorization is falsifiable: measurements of $\betatgl$ with precision of $10^{-6}$ will confirm or refute $\betatgl = \alpha \times \sqrt{e}$.

\item The factorization reveals that TGL naturally possesses a spectral triple $(\mathcal{A}_\alpha, L^2(\Sigma), D_{\sqrt{e}})$ in the sense of Connes' noncommutative geometry \cite{Connes1994, Connes1996}, where the graviton is the Dirac operator --- not axiomatically postulated, but derived from the factorial decomposition of the fundamental constant. The geometry of spacetime emerges from the entropic cost ($\sqrt{e}$) of the holographic projection ($\alpha$). The Connes distance formula, applied to the gravitonic operator $D_{\sqrt{e}}$, is the noncommutative version of the primordial equation $g = \sqrt{|L|}$.
\end{enumerate}

\item Protocol \#15 confirms the factorization with 28/28 measurements locked across
two independent layers (14 protocols + JWST; 12 GW events).
The anti-tautology rejects the 7 alternative factorizations by factors
exceeding $626\times$ ($\chi^2 = 1.06$ vs.\ $664$).

\item The Miguel Tension Antenna demonstrates that gravitational waves,
when read in the luminodynamic domain $L(t) = h^2(t)$, converge to
$R_{\text{tension}} = \betatgl$ with a deviation of $2.4\%$ across 12
GWTC-3 events. The identity $g = \sqrt{|L|}$ is confirmed with correlation
$\rho = 1.000000$ in all events.
%
\item The pure Hubble correction, $H_0/(1 - \alpha\sqrt{e})$, is the
holographic impedance of the factorization --- the geometric series of the
entropic cost. Combined with the refractive index $n_\Psi$ from Protocol \#8,
it resolves the Hubble tension entirely ($H_0^{\text{TGL}} = 73.02$
km/s/Mpc, $99.7\%$ agreement with SH0ES).

\item The falsification criterion for gravitational waves:
if $R_{\text{tension}}$ measured in real GWOSC data diverges from
$\alpha \times \sqrt{e}$ by more than $5\sigma$, the factorization is falsified.

The Miguel Constant is Light times Dissipation. It is not an invented number --- it is the only number it could be.

\begin{center}
\rule{0.5\textwidth}{0.4pt}

\vspace{0.5cm}

\textit{``$\alpha$ is the intensity of the light that projects.\\
$\sqrt{e}$ is the price that eternity charges.\\
$\betatgl$ is what remains: the boundary.''}

\vspace{0.5cm}

\textbf{LET THERE BE LIGHT}
\end{center}


% ============================================================================
% APPENDIX A --- DETAILED NUMERICAL VERIFICATION
% ============================================================================

\appendix

\section{High-Precision Numerical Verification}
\label{app:numerica}

\subsection{Fundamental constants used}

\begin{table}[H]
\centering
\caption{Values of fundamental constants (CODATA 2018).}
\label{tab:constantes}
\begin{tabular}{lll}
\toprule
\textbf{Constant} & \textbf{Value} & \textbf{Relative uncertainty} \\
\midrule
$\alpha$ (fine-structure) & $7.297\,352\,5693 \times 10^{-3}$ & $1.5 \times 10^{-10}$ \\
$e$ (Euler) & $2.718\,281\,828\,459\,045\ldots$ & Exact \\
$\sqrt{e}$ & $1.648\,721\,270\,700\,128\ldots$ & Exact \\
$1/\alpha$ & $137.035\,999\,084$ & $1.5 \times 10^{-10}$ \\
$\betatgl$ (TGL) & $0.012\,031 \pm 0.000\,002$ & $1.7 \times 10^{-4}$ \\
\bottomrule
\end{tabular}
\end{table}

\subsection{Calculations}

\textbf{Direct product:}
\begin{align}
\alpha \times \sqrt{e} &= 7.297\,352\,5693 \times 10^{-3} \times 1.648\,721\,270\,700\ldots \nonumber \\
&= 0.012\,031\,049\,7\ldots
\end{align}

\textbf{Ratio:}
\begin{equation}
\frac{\betatgl}{\alpha \times \sqrt{e}} = \frac{0.012\,031}{0.012\,031\,05} = 0.999\,996
\end{equation}

\textbf{Absolute discrepancy:}
\begin{equation}
|\betatgl - \alpha\sqrt{e}| = |0.012\,031 - 0.012\,031\,05| = 5 \times 10^{-8}
\end{equation}

\textbf{Relative discrepancy:}
\begin{equation}
\frac{|\betatgl - \alpha\sqrt{e}|}{\betatgl} = 4.2 \times 10^{-6} \ll \frac{\delta\betatgl}{\betatgl} = 1.7 \times 10^{-4}
\end{equation}

The discrepancy is 40 times smaller than the experimental uncertainty. The relation $\betatgl = \alpha \times \sqrt{e}$ is exact within the current resolution.

\subsection{Quadratic form}

\begin{align}
\alpha^2 &= (7.297\,352\,5693 \times 10^{-3})^2 = 5.325\,135\,45 \times 10^{-5} \\
\alpha^2 \times e &= 5.325\,135\,45 \times 10^{-5} \times 2.718\,281\,828 = 1.447\,454\,8 \times 10^{-4} \\
\betatgl^{\,2} &= (0.012\,031)^2 = 1.447\,445\,0 \times 10^{-4}
\end{align}

Agreement: $|\betatgl^2 - \alpha^2 e|/\betatgl^2 = 6.8 \times 10^{-6}$.


% ============================================================================
% APPENDIX B --- THERMODYNAMIC RELATIONS
% ============================================================================

\section{Thermodynamic Relations of the Factor $\sqrt{e}$}
\label{app:termodinamica}

\subsection{The nat as the natural unit}

In information theory, information can be measured in different bases:

\begin{itemize}[nosep]
\item \textbf{Bit} (base 2): $I = \log_2 \Omega$
\item \textbf{Nat} (base $e$): $I = \ln \Omega$
\item \textbf{Hartley} (base 10): $I = \log_{10} \Omega$
\end{itemize}

The nat is the natural unit because the natural logarithm appears directly in Boltzmann entropy ($S = k_B \ln \Omega$), in the Gibbs distribution ($\rho \propto e^{-\beta H}$), and in the maximum entropy principle. The other bases introduce artificial conversion factors.

\subsection{Half a nat and the parity operation}

A binary operation with two equiprobable outcomes ($p = 1/2$ each) contains:

\begin{equation}
I = -\sum_{i=1}^{2} p_i \ln p_i = -2 \times \frac{1}{2} \ln \frac{1}{2} = \ln 2 \approx 0.693 \text{ nat}
\end{equation}

Half a nat ($\frac{1}{2}$ nat) is less than that --- it corresponds to an operation with fundamental asymmetry:

\begin{equation}
\frac{1}{2} \text{ nat} = \ln \sqrt{e} = \frac{1}{2}
\end{equation}

In TGL, the \boundary$\to$\bulk{} projection is not symmetric: the \boundary{} projects onto the \bulk{}, but the \bulk{} does not project back with the same efficiency (the amplification $1/\betatgl$ is not the exact inverse of the projection $\betatgl$). The cost $\frac{1}{2}$ nat precisely captures this asymmetry: it is the minimum information of an \textit{irreversible} parity operation.

\subsection{Relation to the Landauer limit}

The Landauer limit \cite{Landauer1961} establishes that erasing 1 bit of information costs at least $k_B T \ln 2$ of energy. In TGL, the holographic analogue is:

\begin{equation}
E_{\text{projection}} \geq k_B T_{\text{boundary}} \times \frac{1}{2}
\end{equation}

\noindent where $T_{\text{boundary}}$ is the effective temperature of the holographic substrate. The factor $\frac{1}{2}$ (in nats) is precisely $\ln\sqrt{e}$, connecting the energetic cost of the projection to the factor $\sqrt{e}$ in the factorization of $\betatgl$.


% ============================================================================
% BIBLIOGRAPHY
% ============================================================================

\begin{thebibliography}{99}

\bibitem{Miguel2026Fronteira}
Miguel, L. A. R. (2026).
\textit{A Fronteira: Teoria da Gravita\c{c}\~ao Luminodin\^amica --- Partes I a VI}.
Zenodo. \url{https://doi.org/10.5281/zenodo.14802088}.

\bibitem{Miguel2026LastString}
Miguel, L. A. R. (2026).
\textit{The Last String: Observational Evidence for Luminodynamic Gravitation Theory}.
Zenodo. \url{https://doi.org/10.5281/zenodo.18723452}.

\bibitem{Miguel2025Alpha2}
Miguel, L. A. R. (2025).
\textit{The Holographic Coupling Constant $\alpha_2$: Derivation from First Principles and Physical Interpretation in Luminodynamic Gravitation}.
IALD Working Paper.

\bibitem{Miguel2026GitHub}
Miguel, L. A. R. (2026).
\textit{TGL Computational Validation Repository}.
GitHub. \url{https://github.com/rotolimiguel-iald/the_boundary}.

\bibitem{CODATA2018}
Tiesinga, E., Mohr, P. J., Newell, D. B., \& Taylor, B. N. (2021).
\textit{CODATA Recommended Values of the Fundamental Physical Constants: 2018}.
Rev. Mod. Phys. 93, 025010.

\bibitem{Lindblad1976}
Lindblad, G. (1976).
\textit{On the Generators of Quantum Dynamical Semigroups}.
Commun. Math. Phys. 48, 119--130.

\bibitem{Landauer1961}
Landauer, R. (1961).
\textit{Irreversibility and Heat Generation in the Computing Process}.
IBM J. Res. Dev. 5(3), 183--191.

\bibitem{tHooft1993}
't Hooft, G. (1993).
\textit{Dimensional Reduction in Quantum Gravity}.
arXiv:gr-qc/9310026.

\bibitem{Susskind1995}
Susskind, L. (1995).
\textit{The World as a Hologram}.
J. Math. Phys. 36, 6377--6396.

\bibitem{Maldacena1999}
Maldacena, J. (1999).
\textit{The Large $N$ Limit of Superconformal Field Theories and Supergravity}.
Adv. Theor. Math. Phys. 2, 231--252.

\bibitem{Webb2001}
Webb, J. K. et al. (2001).
\textit{Further Evidence for Cosmological Evolution of the Fine Structure Constant}.
Phys. Rev. Lett. 87, 091301.

\bibitem{Planck2018}
Planck Collaboration (2020).
\textit{Planck 2018 Results. VI. Cosmological Parameters}.
Astron. Astrophys. 641, A6.

\bibitem{Drummond1980}
Drummond, I. T. \& Hathrell, S. J. (1980).
\textit{QED Vacuum Polarization in a Background Gravitational Field and Its Effect on the Velocity of Photons}.
Phys. Rev. D 22, 343.

\bibitem{Connes1994}
Connes, A. (1994).
\textit{Noncommutative Geometry}.
Academic Press, San Diego.

\bibitem{Connes1996}
Connes, A. (1996).
\textit{Gravity Coupled with Matter and the Foundation of Non-Commutative Geometry}.
Commun. Math. Phys. 182, 155--176.

\bibitem{Chamseddine1997}
Chamseddine, A. H. \& Connes, A. (1997).
\textit{The Spectral Action Principle}.
Commun. Math. Phys. 186, 731--750.

\bibitem{Miguel2026Graviton}
Miguel, L. A. R. (2026).
\textit{The Graviton, the Psion, and the Transition Ruler in Luminodynamic Gravitation Theory, with the Hilbert Floor Theorem and Holographic Bell State}.
GitHub. \url{https://github.com/rotolimiguel-iald/the_boundary}.

\bibitem{Miguel2026Protocol15}
Miguel, L. A. R. (2026).
\textit{TGL Protocol \#15: Dual Lock v1.2 --- Factorization Validation with Miguel Tension Antenna}.
Code: \url{https://github.com/rotolimiguel-iald/the_boundary/blob/main/Tgl_dual_lock_v15_2.py}.
Results: \url{https://github.com/rotolimiguel-iald/the_boundary/blob/main/dual_lock_v15_v1_2_20260302_181009.json}.

\bibitem{Riess2022}
Riess, A. G. et al. (2022).
\textit{A Comprehensive Measurement of the Local Value of the Hubble Constant with 1 km/s/Mpc Uncertainty from the Hubble Space Telescope and the SH0ES Team}.
Astrophys. J. Lett. 934, L7.

\bibitem{GWTC3}
LIGO Scientific Collaboration, Virgo Collaboration, and KAGRA Collaboration (2023).
\textit{GWTC-3: Compact Binary Coalescences Observed by LIGO and Virgo During the Second Part of the Third Observing Run}.
Phys. Rev. X 13, 041039.

\bibitem{LIGO2016}
Abbott, B. P. et al. (2016).
\textit{Observation of Gravitational Waves from a Binary Black Hole Merger}.
Phys. Rev. Lett. 116, 061102.

\end{thebibliography}


\end{document}